\chapter{The neutral gas transport equation; Monte-Carlo terminology} \label{sec1} 
\textbf{General Remarks}
\medskip


To introduce the terminology used throughout this report, we
briefly recall the basic definitions and principles of a Monte
Carlo linear transport model, following the lead of many textbooks
on Monte Carlo methods for computing neutron transport (see e.g.,
\cite{kn:Spanier69a}). The original derivation of the EIRENE Monte Carlo
model, building on these neutronics concepts and terminology,
was carried out in D. Reiter's phd thesis (in german) \cite{kn:Reiter84a}
in the early eighties of the last century.
We begin with the linear transport equation
for the pre-collision density, written as integral equation
(linear non-homogeneous Fredholm integral equation of 2nd kind).
Distinct from standard terminology in the (analytic) transport
theory we do not discuss analytic properties of the various terms
in the equation, but, instead, point out their probabilistic
interpretation, as needed for a Monte-Carlo solution of that
equation. Next (\cref{sec1.2}) we sketch the Monte Carlo
procedure to solve such equations, by referring to the two most
often applied techniques: ``track-length - and collision based
estimators''. In the third subsection we briefly describe the
treatment of boundary conditions (models for interaction of the
particles with surrounding surfaces) and discuss some special
models which are in use for the neutral gas transport in fusion
plasma devices. Then (\cref{sec1.4}) we discuss the most
important source function (non-homogeneous part of the integral
equation) and its implementation in a Monte Carlo algorithm,
namely the surface source of neutral particles due to
recombination of ions incident on solid surfaces at the boundary
or inside of the computational area. In \cref{sec1.5} we
comment on the implementation of geometry within the framework of
a Monte Carlo code in general terms and for the EIRENE code in
particular. In \cref{sec1.6} the time dependent mode of
operation of the EIRENE code is described. It merely amounts just
to an increase in the dimensionality by one, by adding a time
co-ordinate and treating it, formally, in a rather symmetric
fashion with the other spatial coordinates. See also reference
\cite{kn:Reiter94}.
\\
Two kinds of non-linearities may be accounted for:
\\
In \cref{sec1.7} the nonlinear behavior resulting from
background data, which depend on neutral particle transport
(sources and sinks), is discussed. The algorithm of the B2-EIRENE
code system \cite{kn:Reiter92b}
is described. More details on this can be found in the report \cite{kn:Maddison94}.\\
In the final section of this introductory \cref{sec1.8} the
non-linear \ac{BGK} formalism and the \ac{DMCS} method
for self collisions between neutral particles, as implemented in
the EIRENE code, is described.
\newpage

\section{The linear Boltzmann equation for the distribution function
\texorpdfstring{$f$}{f}} \label{sec1.0}
EIRENE solves a multi species set of coupled ``Boltzmann''-type equations, in arbitrary 3D geometries.
Strictly speaking it is the Boltzmann equation generalized from its original single species form with
bi-linear collision kernel for elastic collisions,
to a far more complicated collision integral, which also represents  ``chemical reactions''.
This generalisation is also often referred to as \ac{WCU} equation in literature \cite{kn:WCU}.
In this present section we start with the Boltzmann equation. We then deal with linearisations of collision
kernels by fixing the phase space distribution of one of the two collision partners (which we refer to as ``bulk''
or ``background'' or even simply ``plasma'' species). Generalization to the full \ac{WCU} type multi-species collision kernels, which in the same
linearized form provide the mathematical description of the equations solved by EIRENE, is then discussed next, \cref{sec1.01}.

Relaxation of the linearisation assumptions is needed only when there is self-interaction amongst the species considered by EIRENE,
as it is required e.g. when radiation transfer is included (coupling between ``photons'' and atoms) or if
neutral-neutral collisions are relevant. In this latter case so called \ac{BGK} approximations to the full collisions integral
are employed, and the non-linearity is dealt with by successive linearisation (i.e. by iteration).

The strongest non-linearity in EIRENE applications is typically the back-reaction of the ``plasma-background'' (the host medium) on the
neutral gas and radiation fields. This, however, is dealt with in an operator splitting cycling procedure (see \cref{sec4}) between a plasma solver (``diffusion-advection sub-module" solver, for given reaction terms) and EIRENE (``reaction part'' solver, for given diffusion advection solution, i.e. for given background medium), in which
EIRENE, within each single cycle, may still be operated in linear mode.






The term \textsl{ $\mu$-space } in this report refers to the phase
space of a single particle. The quantity of interest is then the
\textsl{ one particle distribution function }  \textit{f}
\[
f(\vr, \vv, i, t)  \quad \mbox{or} \quad \tilde f(\vr, E, \vOmega, i, t)
\quad \mbox{or} \quad f(x) ,
\]
where the state $x$ of the $\mu$-space is characterized by a
position vector $\vr$, a velocity vector
$\vv$, a species index $i$ ($i$ stands for, e.g.,
H, D, T, D$_2$, DT, He, CH$_n$,...) and the time $t$.

The number density $n_i(\vr)$ for species $i$ then
reads:
\[
n_i(\vr,t) = \int d^3{v} \ f(\vr,\vv,i,t)
\]
\\
Instead of $\vv$ we sometimes utilize the kinetic energy
$E$, $E=m/2 \cdot v^2$, and the unit (speed) vector
$\vOmega = \vv/|\vv|$ in the
direction of particle motion. Hence:
\begin{equation}\nonumber
f(\vr)d^3v = \tilde f(E, \vOmega)dEd\Omega \quad \mbox{where}  ~~d^3v = v^2dv d\Omega \quad \mbox{and} ~~ d\Omega = \sin \theta d\theta d\phi
\end{equation}
and
\[
f(\vr, \vv, i ,t)   = \left(\frac{m}{v}\right)
\tilde{f}(\vr, E, \vOmega, i, t)
\]
The distribution function in the form of $\tilde f$ clearly remains meaningful also in the case of massless particles (photons),
i.e. for applications of EIRENE to radiation transfer problems.
\\
The notation used now may appear to be unnecessarily complicated, but we indent to prepare here already the
more general discussion in \cref{sec1.01} on the \ac{WCU} equation.

We start with the original (elastic collisions only) ``Boltzmann Equation" \cite{kn:Cercignani},
but by assuming additionally that collisions are discontinuous events (i.e.: finite range
interaction potentials, or, at least, that proper cut-off procedures are applicable).
This additional assumption allows us to separate, in the Boltzmann collision integral
given below, the collisional ``loss" and ``gain" terms into two
separate integrals. Otherwise cases could arise in which only the net collision term $\int\!\!\int\!\!\int \dots (f'f_b'-ff_b)$ would lead to finite results,
but not the collision term in the form as given in \cref{eq1a} below: $\int\!\!\int\!\!\int \dots (f'f_b')- \int\!\!\int\!\!\int \dots (ff_b)$.

Further: Lets consider only one specific species $i_0$, now sometimes
omitting this species index. We assume that there are only
collisions of this species $i_0$ with only one further species, here
labeled $b$. Generally $b=i_0$ is also possible, but anticipating that later, when discussing \textbf{linear} Boltzmann models,
these species $b$ will be referred to as background species, plasma species, etc.., with an externally given (fixed) distribution $f_b(\vx,\vv,t)$.
Restriction, for the time being, to elastic collisions means that exactly one particle of each of these two
species $i_0, b$ will also be present after the collision event (i.e.,
chemical reactions, inelastic collisions with change of internal energy,
are excluded for the time being,
but their description within a ``Boltzmann-like" framework is discussed below,
\cref{sec1.01}) on the so called ``\ac{WCU}" equation.

The familiar Boltzmann
equation for the distribution function $f=f_{i_0}$ for a particular species (a particular component of a gas mixture) $i_0$
reads

% gleichungsnummern fest auf eq. 1a setzen
\renewcommand{\theequation}{1.1a}
\begin{align} \label{eq1a}
\lbrack\frac\partial{\partial t} + \vv \cdot
\nabla_{\vr} +
\frac{\vF(\vr,\vv,t)}{m} \cdot
\nabla_{\vv}\rbrack \ f(\vr,\vv,t)
& = 
\int\!\int\!\int
\sigma(\vv',\vV';\vv,\vV)
|\vv'-\vV'| f(\vv')
f_b(\vV')  \nonumber \\
& -  \int\!\int\!\int
\sigma(\vv,\vV;\vv',\vV')
|\vv-\vV| f(\vv) f_b(\vV)
\nonumber \\  
& + Q(\vr,\vv,t)
\end{align}

% gleichungsnummern einen hochzaehlen.
%\addtocounter{equation}{+1}
%  gleichungszaehler auf default mode zuruecksetzen
\renewcommand{\theequation}{\arabic{section}.\arabic{equation}}
There is one such equation for each ``particle species $i$" considered, but for elastic collision
without exchange between species. Always only two of them (here: $i_0$ and $b$) a directly coupled to each other.


$Q \ (= Q_{i_0})$ is any external source (particles of species $i_0$ injected per unit volume  in phase space and per unit time).
\

Integration in the first term on the right hand side (``gain term") of this equation is over the (pre-collision) velocities $\vv',
\vV'$  as well as over one of the two post collision velocities: $\vV$, the velocity of species $b$.
We write more suggestively  $\sigma(\vv',\vV'; \vv,\vV) = \sigma(\vv',\vV' \rightarrow \vv,\vV)$
for the (multi-) differential cross section for a binary particle collision process. The product of this $\sigma$ with the relative pre-collision velocity
$|\vv'-\vV'|$ is the transitional (collision) probability leading from velocities prior to the collision $(\vv',\vV')$ to the pair $(\vv,\vV)$ of velocities after the collision.

In a multi-component gas mixture the first two arguments in $\sigma$, namely the velocities
$\vv',\vV'$ in the first integral, correspond to the species $i_0$
and $b$, respectively. Via elastic collisions these are turned into
the post collision velocities $\vv, \vV$,
again for the same species $i_0$ and $b$, respectively. The first integral,
therefore, describes the total rate of collisional transitions
into the velocity space interval
$[\vv,\vv+d\vv]$ for species $i_0$

In the second integral on the right hand side (``loss term") the role of pre and post collision velocities is exchanged, but the integration
remains over the same three velocities $\vv',\vV'$ and $\vV$, i.e. integration is now over all post collision velocities and
over all incident velocities $\vV$ of collision partner $b$. Hence this
integral describes the total collisional rate of loss for species $i_0$ from that phase space interval $[\vv,\vv+d\vv]$.
Furthermore, $m = m_{i_0}$ is the particle mass and
$\vF(\vr,\vv,t)$ is a volume force
field. The fact that the cross section $\sigma$ remains the same in the forward (gain) and reverse (loss) integral
after exchanging the pre-collision with the post collision velocity arguments
follows from time-reversal (detailed balancing) in case of elastic collisions.


The four velocity arguments in both collision integral terms are not truly independent:
the conservation laws for total energy and
momentum in a collision event must be fulfilled.
\begin{align}\label{massenergy}
m_{i_0} \vv' + m_b\vV' &= m_{i_0}\vv + m_b \vV \nonumber\\
\frac12 m_{i_0}v'^2+\frac12m_bV'^2&=\frac12 m_{i_0}v^2 + m_b V^2 +\Delta_E
\end{align}
Here $\Delta_E$ is the exchange of internal energy in the collision (on expense of the energy of relative motion), and, by definition, $\Delta_E = 0$ for elastic collisions.
We may assume that cross section $\sigma$ contains appropriate delta function factors such that the three
integrations over velocity spaces reduce to integrations over lower dimensional manifolds on which these conservation laws are fulfilled.
The modulus of the relative velocity $g$ is unchanged in elastic collisions,
$g_{i_0,b} = |\vv-\vV| = |\vv'-\vV'| = g_{i_0,b}'$,
and hence the collision is just a rotation of this relative velocity, by an angle $\vec\Omega$. 6-fold integration over full velocity spaces $\vv', \vV'$
in $\Re^6$ reduce to just the 2 fold integration
over $\Omega:  d\Omega=\sin\chi d\chi d\epsilon$ with $\chi$ being the deflection angle
and $\epsilon$ the scattering angle.


Writing
\begin{equation} \nonumber
D_t^{i_0} = \frac\partial{\partial t} + \vv \cdot
\nabla_{\vr} + \frac{\vF_{i_0}(\vr,\vv,t)}{m_{i_0}} \cdot
\nabla_{\vv}
\end{equation}
for the left hand side (convective part) of \cref{eq1a}, leads to the Boltzmann equation in more commonly used text book form:

% gleichungsnummern fest auf eq. 1.1a1 setzen
\renewcommand{\theequation}{1.1a1}
\begin{align} \label{eq1a1}\nonumber
D_t^{i_0} \ f_{i_0}(\vr,\vv,t)
& = 
\int_{\Re^3,\Omega}  d\vV'd\Omega \ \
\sigma(\vv,\vV;\vv',\vV')
|\vv-\vV| \left[f_{i_0}(\vv')
f_b(\vV')- f_{i_0}(\vv) f_b(\vV)\right]
\nonumber \\  
& + Q_{i_0}(\vr,\vv,t)
\end{align}

% gleichungsnummern einen hochzaehlen.
%\addtocounter{equation}{+1}
%  gleichungszaehler auf default mode zuruecksetzen
\renewcommand{\theequation}{\arabic{section}.\arabic{equation}}

and the velocities $\vv',\vV $ are uniquely determined by $\vv, \vV'$ and scattering angles $\Omega$ via mass and energy conservation \cref{massenergy}.


The integral in this equation, or, equivalently, the first two terms on the right hand side of \cref{eq1a}, comprise the Boltzmann collision integral $\frac{\delta
f_{i_0}}{\delta t}|_{_{coll}}$. If there are more than just one possible type
of collision partners (species) ``$b$", then the collision integral has to be
replaced by a sum of collision integrals over all collision
partners $b$, including, possibly, $b=i_0$ (self collisions)
\begin{equation}\label{collop}
\frac{\delta f_{i_0}}{\delta t}|_{_{coll}}= \frac{\delta f_{i_0}}{\delta t}|_{_{gain}}- \frac{\delta f_{i_0}}{\delta t}|_{_{loss}}= \sum_b \frac{\delta f_{i_0}}{\delta
t}|_{_{coll_b}}.\end{equation}

All these collision operators are bi-linear in the distribution
functions $f_{i_0}, f_b$. The first term on the right hand side is due to
scattering into the element $d\vv$ of velocity space and
we shall abbreviate it by defining the collision kernel
(``redistribution function") $C_{i_0}$ for $i_0$ particles alone, i.e. with an operator acting on (pre-collision) $f_{i_0}$ alone, already containing the proper integrations over pre-
and post collision velocities $\vV', \vV$ of species $b$ collision partners:
\begin{equation}\label{colkern}
\frac{\delta f_{i_0}(\vv)}{\delta t}|_{_{gain}}
 =
 \int d^3{v}'  C_{i_0}(\vv' \rightarrow \vv)
|\vv'| f_{i_0}(\vv')
 =
\int d^3{v}' \sum_b \ C_{i_0,b}(\vv' \rightarrow \vv)
|\vv'| f_{i_0}(\vv')
\end{equation}
In more general situations (including chemical reactions) in addition to being an integral over particle $``b"$ velocity distributions, the kernels $C$ can be quite complicated integrals due to collision kinetics, as it then involves
not only multiple differential cross sections, but also, possibly,
particle multiplication factors, e.g. in case of fission by
neutron impact, dissociation of molecules by electron impact, or
stimulated photon emission from excited atoms. It can also
include absorption, in which case the post collision state must be
an extra ``limbo" state outside the phase-space considered. Due to
both particle multiplication and/or absorption the collision
kernel $C$ is not normalized to one, generally.

The second term on the right hand side of the Boltzmann equation $\frac{\delta f}{\delta t}|_{_{loss}}$ is, formally, much simpler, because
the function $f_{i_0}(\vv)$ can be taken out of the integrals.
It is often convenient to even take the product $|\vv| \cdot f_{i_0}$ before the
integral. The remaining integral is then just the total
``macroscopic cross section" $\Sigma_t$ (dimension: 1/length), i.e., the inverse local mean
free path. Apart from dependence on individual velocity $\vv$  of particle $i_0$ it is solely defined by total
cross sections averaged over velocity distributions of particles ``$b$", i.e. by ``reaction rates", and is independent of particle multiplication factors, since we
only consider binary collisions (exactly two pre-collision partners always, which lead to the loss of exactly
one particle of species $i_0$ from phase space).

This term is then, often, taken on the left hand side of the
Boltzmann equation with a positive sign, in the form:
\begin{equation}\label{extcoef}
\frac{\delta f_{i_0}}{\delta t}|_{_{loss}}
 =
\Sigma_{t,i_0}(\vr,\vv) |\vv|
f_{i_0}(\vv)
\end{equation}

Note, however, that in case of ``self collisions" within the community of particles $i_0$, the macroscopic cross section $\Sigma_t$
itself is a (nonlinear) function of the dependent
variable $f_{i_0}: \Sigma_{t,i_0}(\vr,\vv,f_{i_0})$,
since it then contains also an integration (over velocities $\vV$ of the collision partners) of $\sigma(...,\vV \rightarrow ...) \times f_{i_0}(\vV)$.

\paragraph{The linear form of Boltzmann transport equation}\hfill

\vspace{1em}

If all distributions $f_b(\vV)$ of  the collision partners $b$ of species $i_0$ are assumed to be given, then the
kernel $C = C_{i_0}$ defined in \cref{colkern} is linear and the expression given there defines  a linear
integral (``collision") operator acting on the distribution $f_{i_0}(\vv)$.
For the linear case, i.e. for $f_b$ given for all collision partners other than $i_0$, and for self collisions ($i_0 + i_0 \rightarrow ...$) being excluded, also the ``extinction coefficient"
$\Sigma_t = \Sigma_{t,i_0}$ defined by \cref{extcoef} is independent of the dependent variable $f = f_{i_0}$,
and this term (out-scattering) just describes the loss of particle
flux of $i_0$ particles due to any kind of interaction of them with any species $b$ from
the background ``host" medium.

With these formal substitutions the Boltzmann equation for distribution function $f(\vv) = f_{i_0}(\vv)$ takes a
form which is often more convenient, in particular in linear
transport theory (from now on omitting subscript $i_0$ on all terms):

% gleichungsnummern fest auf eq. 1.1b setzen
\renewcommand{\theequation}{1.1b}
\begin{align} \label{eq1b}
\lbrack\frac\partial{\partial t} + \vv \cdot
\nabla_{\vr} +
\frac{\vF(\vr,\vv,t)}{m} \cdot
\nabla_{\vv}\rbrack \ f(\vr,\vv,t) +
\Sigma_t(\vr,\vv) |\vv|
f(\vv)
  &=
\int d^3{v}' \ C(\vv' \rightarrow \vv)
|\vv'|
f(\vv') \nonumber \\
  &+ Q(\vr,\vv,t)
\end{align}
%  gleichungszaehler auf default mode zuruecksetzen
\renewcommand{\theequation}{\arabic{section}.\arabic{equation}}
% gleichungsnummern einen runterzaehlen.
\addtocounter{equation}{-1}
which, in the linear case (no self collisions, $f_b$ given externally) is a linear equation for the distribution
function $f = f_{i_0}(\vx,\vv,t)$ for species $i_0$.
\\
Clearly a computationally crucial simplification is provided in this linear case, which means neglect of all
interactions within the community of species $``i"$, retaining only: $i_0 + b \rightarrow \cdots$ events.
In practice for neutral particle (species group $i$) transport in plasmas (species group $b$) this mostly means neglect of neutral-neutral interaction,
retaining only neutral-plasma collisions (for given plasma conditions).
For the status of options in EIRENE to deal with ``non-linear''
``self-collisions" ($i_0 + i_0 \rightarrow \cdots$)  or with ``cross collisions" ($i_0 + i_1 \rightarrow \cdots$) amongst the
species for which the kinetic equation is solved, see below, \cref{sec1.8}. Note that also combined (consistent) neutral particle and
radiation transport models fall into the class of non-linear Monte Carlo applications, except for optically thin media in which case the photon gas
and the neutral particle gas are decoupled.

\paragraph{stationary kinetic transport equation}\hfill

\vspace{1em}

Often the characteristic time constants for neutral particle transport
phenomena are very short (\si{\micro\second}), compared to those for plasma
transport (ms). We can, therefore, often neglect explicit time
dependence in the equations describing the neutral particles. This
is done in most applications. The extension to time-dependent
problems is rather straight forward and the procedure in the
EIRENE code for such cases is described in \cref{sec1.6} below.
This is because in a Monte Carlo transport procedure
time can be treated in a completely symmetric fashion with the spatial coordinates.
\\
For stationary (time-independent) problems the (scalar) transport flux (``angular flux") $\Phi$, where
\begin{subequations}
\begin{equation}\label{trans-flux}
\Phi (x)   :=   | \vv |  \cdot f( \vr , \vv , i)   ,
\end{equation}
with, again, $| \vv| \Omega = \vv$,
is sometimes used as dependent variable, in preference to the distribution $f(x)$.
For example the \emph{total particle flux} is obtained from this ``angular flux'' by integrating it over all velocity space angles $\Omega$,
i.e. as the zeroth $\Omega$ - moment of $\Phi(\vr,\vv,i)$.

In particular for stationary ($\partial/\partial t = 0$)  and force free ($\vF = 0$) problems, as e.g. often encountered in linear transport theory such as
neutronics, radiation transport, neutral particle transport in plasmas, etc.,
the transport equation then reduces to the more compact form:
% gleichungsnummern fest auf eq. 1.1c setzen
\renewcommand{\theequation}{1.1c}
\begin{equation} \label{eq1c}
\nabla_{\vr}
 \Phi(\vr,\vv,t) +
\Sigma_t(\vr,\vv)
\Phi(\vr,\vv)
  =  Q(\vr,\vv,t) +
\int d^3{v}' \ C(\vr;\vv' \rightarrow \vv)
\Phi(\vr,\vv')
\end{equation}
%  gleichungszaehler auf default mode zuruecksetzen
\renewcommand{\theequation}{\arabic{section}.\arabic{equation}}
% gleichungsnummern einen runterzaehlen.
\addtocounter{equation}{-1}

Alternatively, in computational domains with non-vanishing
collisionality (i.e., if $\Sigma_t(x) \ne 0$ everywhere) the (pre-)
collision density $\Psi$ is used, i.e.,
\begin{equation}\label{defPsi}
\Psi (x)
= \Sigma_{t} (x)  \cdot \Phi (x) =\nu_t(x) \cdot f(x)  , \hspace{1cm} \Sigma_t =
\Sigma_t(\vr,E,\vOmega,i)= \nu_t(\vr,E,\vOmega,i)/|\vv|
\end{equation}
\end{subequations}
where, again, the ``macroscopic cross section" $\Sigma_t$ is the
total inverse local mean free path (dimension: 1/length), and $\nu_t$ is the collision frequency (dimension: 1/time). This
cross section can be written as a sum $ \Sigma_t   =   \sum
\Sigma_k$ over macroscopic cross sections for the different types
(identified by the index $k$) of collision processes. Further
details about this ``macroscopic cross section" and its relation
to the conventional (``microscopic") cross sections are given
below, \cref{sec1.2.3}.

In closing this section we note:


All these equations \cref{eq1a,eq1b,eq1c,eq1d}
given below are equivalent, of course,
so are their corresponding \ac{WCU} generalizations discussed in the next \cref{sec1.01}.
Which particular
form is used in a particular discussion depends upon the issue which is
considered. For example, the collision estimator for evaluating moments (responses) of the solution to the Boltzmann equation can be shown to be
unbiased quite conveniently when using eq. \cref{eq1d}, whereas the track-length
estimator, used for the same purpose in EIRENE, is most easily understood with the from \cref{eq1d}. Time dependent
cases are best discussed utilizing the form \cref{eq1b}, etc...

\subsection{The \acs{WCU} generalisation of the Boltzmann equation} \label{sec1.01}

The mathematical generalization from the classical Boltzmann equation
(for a system undergoing only elastic collisions),
equation \cref{eq1a}
to the semi-classical Boltzmann
equation (the \ac{WCU} equation, \cite{kn:WCU}) for accomodating also chemical reactions (including, for example,
vibrational relaxation or exchange of internal energy $\Delta E$ as special
cases, but also chemical reactions $A+B \rightarrow C+D$, etc.)
is obtained by a proper generalization of the species labels.
The transitions are now symbolized as $\{i'l'\},\{b'l_1'\} \rightleftarrows \{il\},\{bl_1\}$. These
generalized species indices label both the chemical species ($i,b$) and/or the
internal quantum state ($l,l_1$) corresponding to modes of vibration, rotation or excitation.
I.e. we regard two particles (objects) as different, with a different label,
if either they belong to a different
chemical species or they differ by their internal (electronic, or ro-vibrational, etc) state,
as appropriate.


Now, by analogy to the discussion of the original Boltzmann equation \cref{eq1a}, \cref{eq1a1},
we consider then the phase-space balance equation for a given species/internal mode $\{i_0l_0\}$. Then the sum in the collision
integral is over all collision labels $\{i'l'\}, \{b'l_1'\}, \{bl_{1}\}$ and over all velocities of these three ``species".
Even more generally:
in case there are more than two post collision particles, i.e. more than two particles in the exit channel of the collision, integration (velocities) and summation (generalized species indices) in the \ac{WCU} collision integral is over all involved particles, except over $\{i_0l_0; \vv\}$, the parameters of the dependent variable $f_{i_0}(l_0,\vv)$. Employing the notations as in \cref{eq1a1} the full \ac{WCU} prototypical system of kinetic transport equations then reads (written here for two post collision particles only)

\begin{align}\label{WCU}
D_t^{i_0l_0} \ f_{i_0l_0}(\vr,\vv,t)
& = 
\sum_{i',l',b,l_1,b',l_1'}\int_{\Re^3,\Omega}  d\vV'd\Omega \ \
B_{i_0l_0,bl_1}^{i'l',b'l_1'}(\vv,\vV) \ \times \\
& \phantom{=} \times
\left[f_{i'l'}(\vv')
f_{b'l_1'}(\vV')- f_{i_0l_0}(\vv) f_{b,l_1}(\vV)\right]
\nonumber \\
& + Q_{i_0l_0}(\vr,\vv,t) \nonumber
\end{align}
with
\begin{equation}
B_{il,bl_1}^{i'l',b'l_1'}(\vv,\vV)=\sigma_{il,bl_1}^{i'l',b'l_1'}(\vv,\vV;\vv',\vV') \
|\vv-\vV|
\end{equation}

The cross sections in the
corresponding collision integrals
$\sigma_{il,bl_1}^{i'l',b'l_1'}(\vv',\vV',\vv,\vV)$
are multiple differential for scattering at a certain solid angle and post
collision energies with simultaneous transition from $(il, bl_1)$
to $(i'l',b'l_1')$.
Again, as already in the Boltzmann equation restricted to elastic collisions, the pre and post collision velocities
$\vv',\vV'$ and $\vv,\vV$, respectively, are not independent, but constrained by momentum and energy conservation, and the corresponding constrains
are contained in the multi-differential cross sections, e.g. in form of delta functions.
\\
For more than two post collision objects (e.g. fission, dissociative excitation, stimulated radiation emission, etc...), this notation is readily generalized by
adding more superscript labels and more post collision velocities.



The extra discrete label $l$ introduced here compared to the Boltzmann equation may either be regarded as species index:
$i_0 \rightarrow i_0l_0$ or as other additional discrete independent variable (e.g. polarization, in case of radiation transport):
$f_{il}(\vv) \rightarrow f_{i}(l,\vv)$, whichever is more convenient.

As noted above, further generalizations to include particle
splitting, absorption or fragmentation into more than two post
collision products are straight forward, but can more conveniently
be formulated in the $C$-collision kernel formulation introduced above and used below to relate the transport equation
to a Markovian stochastic (discrete time) process.

\paragraph {linear form of \ac{WCU} equation: the generic EIRENE equation}\hfill

\vspace{1em}

As above, linearisation will be again be obtained by grouping generalized labels $\{il\}$ and $\{bl_1\}$ into two disjunct classes,
and allowing for reactions only between initially (prior to the collision) one member $\{i_0l_0\}$ from class $``\{il\}"$, the community of ``test particles''
and one member $\{b_0l_{1_0}\}$ from the disjunct class $``\{bl_1\}"$, the ``bulk", ``background" medium, or in our case, simply ``plasma" species.

\section{The linear integral equation for the collision density
\texorpdfstring{$\Psi$}{psi}} \label{sec1.1}

By formally integrating the characteristics for  \cref{eq1c} the
same transport equation can also be written in integral form. This formal procedure is outlined below in
\cref{sec_green}.

The resulting integral equation is often most conveniently written for the (ingoing-) collision density $\Psi$ \cref{defPsi} rather than for
transport flux $\Phi$:
% gleichungsnummern fest auf eq. 1.1d setzen
\renewcommand{\theequation}{1.1d}
\begin{equation} \label{eq1d}
\Psi(x) = S(x) + \int dx ' \Psi (x ') \cdot K(x '\rightarrow x)  .
\end{equation}
% gleichungsnummern einen runterzaehlen.
\addtocounter{equation}{-1}
%  gleichungszaehler auf default mode zuruecksetzen
\renewcommand{\theequation}{\arabic{section}.\arabic{equation}}
This equation has the general form of the backward integral
equation of a Markovian jump-process and it is therefore
particularly well suited for a Monte Carlo method of solution. The
formal relation between the integro-differential form \cref{eq1c}
and this integral form is very useful to generalize the Monte
Carlo procedure, e.g., to time-dependent equations, and to
Boltzmann-Fokker-Planck equations (which contain diffusive
contributions or diffusive approximations for some processes, in
addition to the jump processes described by the Boltzmann
collision integral). It also allows to make connection to the so
called ``Green's-functions Monte Carlo'' concept (originally
developed for quantum mechanical problems involving solutions to
the Schr\"odinger equation). A corresponding discussion is
postponed to \cref{sec_green}. A direct intuitive
interpretation of the integral equation is already sufficient to
understand the Monte Carlo method of solution and shall be given
first.

In \cref{eq1d} $x'$ and $x$ are the states at two successive
collisions (jumps). The integral $\int dx'$ in \cref{eq1d} is to be understood as
an integral over all initial coordinates, i.e. over the entire physical space, the full velocity space and a
summation over all species indices. The transition kernel $K$ is
usually decomposed, in our context, into a collision- and a
transport kernel, i.e., $C$ and $T$ , where
\begin{equation} \label{eq1.2}
K(\vr  ' , \vv  ' , i ' \rightarrow
\vr  , \vv , i)   =   C( \vr ' ;
\vv ' , i '   \rightarrow   \vv , i)  \cdot T(
\vv , i; \vr ' \rightarrow \vr )   .
\end{equation}
The kernel $C$ is (excluding normalization) the conditional
distribution for new co-ordinates ($\vv $,i) given that
a particle of species $i'$ and with velocity $\vv'$ has
undergone a collision at position $\vr ' $. This kernel
can further be decomposed into:
\begin{equation} \label{eq1.3}
C( \vr ' , \vv ' , i '   \rightarrow
\vv , i) =   \sum_{k} {p_k  \ C_k ( \vr ' ;
              \vv ' , i '   \rightarrow \vv, i)} \
, \ p_k   =    \frac{\Sigma_k}{\Sigma_t}
\end{equation}
with summation over the index $k$ for the different types of
collision processes under consideration and $p_k$ defined as the
(conditional) probability for a collision to be of type $k$. The
normalizing factor
\begin{equation} \label{eq1.4}
c_k (x ')   =
  \sum_{i}   \int d^3v \  C_k (
\vr ' , \vv ' , i ' \rightarrow
\vv , i ) \   , \  {\cal{C}}_k   =    \frac{1}{c_k}
C_k
\end{equation}
gives the mean number of secondaries for this collision process.
The function $ {\cal{C}}_k$ then is a conditional probability
density. The particle absorption process can conveniently be
described by adding an ``absorbing state" $x_a$ to the $\mu$-space
(generally referred to as ``one-point compactification of this
space" in the language of mathematical topology). This ``limbo
state", once it is reached, is never left again if the kernels $T$
or $C$ are employed as transition probabilities. Formally, an
additional collision kernel $C_a (x \rightarrow x_a )$ and an
absorption probability $ p_a   =   \Sigma_a / \Sigma_t$ must be
included in the collision kernel. The quantity $ \Sigma_a$
comprises all collision processes with no next generation
particles within the community of particles considered by the coupled set of kinetic transport equations.
(Ionisation of a neutral atom is a loss, if the resulting ion is not considered further, or dealt with
by a \ac{CFD} code outside the Monte Carlo procedure).
\\The kernel $T$ describes the motion of the test
particles between the collision events. Let, again,  $
\vOmega'$ denote the unit vector in the direction of
particle flight $\vv'/|\vv'|$, and let $ \vOmega_2'$ and $
\vOmega_3'$ be two further unit vectors such that
these three vectors  form an orthonormal basis at the point
$\vr  '$. Neither velocity nor species change along the transition described by $T$, i.e. $\vv'=\vv$ and $i'-i$.
Omitting the corresponding delta functions in velocity space and the Kronecker delta $\delta_{i'i}$
the transport kernel $T$ then reads as
follows:
\begin{subequations}
\begin{align} \label{eq1.5}
(T(l)=)\ \ T ( \vv' , i'; \vr ' \rightarrow  \vr
) &= \Sigma_t ( \vr  ,\vv', i') \cdot \exp
\left [   - \int_{0}^{l=\vOmega' ( \vr -
\vr ' )} ds \Sigma_t ( \vr '  +  s
\vOmega', \vv', i' )  \right ]
\nonumber\\
  &\phantom{=} \cdot \delta (\vOmega_2' ( \vr  -
\vr ' ))    \cdot \delta ( \vOmega_3'   (
\vr - \vr ' ))
  \cdot   H (\vOmega'
( \vr  - \vr  ' )) \\
 &= \Sigma_t ( \vr  ,\vv' , i') \cdot F( \vv' , i'; \vr ' \rightarrow  \vr
) \quad 0\le l \le \infty \label{2.8b}
\end{align}
\end{subequations}
with $H(x) = 0 $ if $ x < 0 $, and $H(x) = 1 $ if $ x \ge 0$, the
Heaviside step function (the unit step function). The two remaining delta functions restrict the motion to a path in the direction
of the initial velocity $\vv'$.

Thus, although strictly being a conditional (on $x'$) distribution in phase space, for an infinite medium $T$ can be interpreted as the
distribution density for the distance $l$ for a free flight
starting from $\vr '$ to the next point of collision
$\vr  = \vr ' + l \cdot \vOmega'$. We shall frequently omit the arguments $\vv',i'$ to simplify notation, because neither initial velocity
nor species change during a free flight.

For a finite medium this distribution can be generalized to (writing shorter $l$ for $(\vr'+l\vOmega',\vv',i')$:
\begin{equation} \label{eq1.5b}
T(l) = \left\{
\begin{array}{lcr}
\Sigma_t ( l) & \cdot \exp
\left [   - \int_{0}^{l} ds \Sigma_t ( s
 )  \right ],&  l< l_{max} \\
\delta(l-l_{max}) & \cdot \exp
\left [   - \int_{0}^{l_{max}} ds \Sigma_t (   s
 )  \right ],  & l\ge l_{max}
\end{array}
\right.
\end{equation}
Here $l_{max}$ denotes the distance along the flight direction from $\vr'$ to the boundary for the computational domain,
to any internal surface at which the test flight shall be stopped, e.g. for scoring surface fluxes there,  or even to the cell boundary.
The ``Transport Kernel" $T$ has dimension: [1/Dimension of phase space] ($T(x'\rightarrow x)dx$ is a probability).
The function $F$ (Dimension: [length times Dimension of $T$])
defined in expression \cref{2.8b}, and analogously from \cref{eq1.5b}
for a finite medium,  will turn out to be the relevant Green's function
for the transport problem, see \cref{sec_green}.
The integral
\begin{displaymath}
\alpha(\vr',\vr) =
\int_{0}^{\vOmega ( \vr - \vr' )}
ds \Sigma_t ( \vr '  +  s \vOmega )
\end{displaymath}
in equation \cref{eq1.5} is well known as ``optical thickness of
the medium" in linear transport theory.
\\ The inhomogeneity $S$ in equation \cref{eq1d} is, excluding
normalization, the distribution density of first collisions,
whereas the integral term in equation \cref{eq1d} describes the
contribution to $\Psi$ from all higher generations. The quantity $S$
can be written as:
%\begin{subequations}  subeqs. functioniert nicht bei latex2html
% gleichungsnummern fest auf eq. 1.6a setzen
\renewcommand{\theequation}{1.6a}
\begin{equation} \label{eq1.6a}
S(x)  =  \int dx '   Q(x ' )  \cdot T(x ' \rightarrow x) ,
\end{equation}
% gleichungsnummern einen runterzaehlen.
\addtocounter{equation}{-1}
%  gleichungszaehler auf default mode zuruecksetzen
\renewcommand{\theequation}{\arabic{section}.\arabic{equation}}
with a source density $Q$. As the problem is linear, $Q$ can be
normalized to 1 and, thus, $Q$ can be considered a distribution
density in $\mu$ -space for the ``\textsl{primary}" birth points of
particles, as, e.g., opposed to the ``\textsl{secondary}" birth point
distribution (or ``post collision density") $\chi$ of particles
after a collision event

% gleichungsnummern fest auf eq. 1.6b setzen
\renewcommand{\theequation}{1.6b}
\begin{equation} \label{eq1.6b}
\chi (x)   =   \int dx '
\Psi(x ' )  \cdot C(x ' \rightarrow x).
\end{equation}
% gleichungsnummern einen runterzaehlen.
\addtocounter{equation}{-1}
%  gleichungszaehler auf default mode zuruecksetzen
\renewcommand{\theequation}{\arabic{section}.\arabic{equation}}
%\end{subequations}  subeqs. functioniert nicht bei latex2html

It can be shown that a unique solution $\Psi(x)$ exists subject to
appropriate boundary conditions and under only mild restrictions
(basically on the constants $c_k$ and $ p_a $) to ensure that the
particle generation process stays sub-critical.
\\ Usually, a detailed knowledge of $\Phi$ or $\Psi$
is not required, but only a set of ``responses", R, defined by
\begin{equation} \label{eq1.7}
R   = \quad  < \Psi | g_c > \quad  =   \int dx   \Psi (x)
 \cdot g_c (x)   \quad   \left (    = \quad  < \Phi | g_t > \quad  =
\int dx   \Phi (x)  \cdot g_t (x)   \right ),
\end{equation}
where $g_c (x)$, $g_t (x)$ are given ``detector functions".

For example all terms in the plasma fluid equations resulting from
neutral plasma interaction can be written in this way. This can be
seen by considering a numerical grid, composed of M mesh cells (spatial and/or temporal),
for the numerical solution of the fluid equations. The detector functions for many responses
needed for fusion plasma applications are hard wired in EIRENE, generalization to any arbitrary
response by resorting to ``user defined detector functions" is described in \cref{sec3.2.1} for tracklength estimates,
and in \cref{sec3.2.2} for collision estimates.

Lets therefore define an entire set of detector functions $g_m$, one for each mesh cell of an external
code,
each including a characteristic function
\begin{subequations}\label{eq2.10}
\begin{equation}
\label{eq2.10a}
g_m=g \times ch_m(\vr,t)  , \quad m=1,2,...,M, \end{equation}
i.e., $ch_m (\vr,t)=1$
inside the numerical mesh cell (or time interval) labeled with the cell index $m$,
and $ ch_m (\vr,t)  =0 $ outside this cell.
Thus profiles of cell volume averaged responses are readily obtained in a single Monte Carlo run.

Estimates of surface fluxes, or point estimates e.g. in time,
are also included in this concept if proper use is made of delta functions to
reduce dimensionality of the response:
\begin{equation}
\label{eq2.10b}
g_{m,\alpha}=g \times ch_m(\vr,t) \times \delta^\alpha(\vr,t) , \quad m=1,2,...,M,
\end{equation}
Here, e.g.,  a surface cell $m$ would discretise  a surface $S$ characterized by a surface delta function $\delta^S$.
\end{subequations}

Equation \cref{eq1.7} shows that Monte Carlo estimates (tallies) fall into the category of \emph{extensive} quantities
(number of particles, total energy, momentum, etc., per cell). \emph{Intensive }quantities  (flow velocity, temperature, etc...)
obtained by dividing two extensive quantities are typically difficult estimate with Monte Carlo methods
(see: ``ratio estimates", correlation between nominator and denominator,...).
If, however, the extensive quantity in the denominator is known exactly, such as e.g. the cell volume in a computational mesh, then, of course,
obtaining an intensive quantity from an extensive quantity with Monte Carlo is a trivial re-scaling.

Monte Carlo estimates of the (intensive) volumetric source terms in the fluid
equations due to the trace particles (neutral particles, but also
trace ions described by the Monte Carlo procedure) are such
cell averages, surface averages, time averages or point estimates (e.g. in time),
averaged over a sub-manifold with reduced dimensionality. (This works, of course, only if the probability
for particle histories crossing this sub-manifold is not zero).
Surface averages in EIRENE are described in \cref{sec3.2.4}, point averages in time
are described in \cref{sec3.2.3}.
This option to reduce dimensionality of responses usually excludes point estimates in real space,
for which special estimators not mentioned here would be required.

Nevertheless, in concluding, depending upon the numerical algorithm in the
fluid code, one may then have to interpret these Monte Carlo estimates properly, e.g.
they may have to be interpolated to the grid points, or be properly rescaled in case of time-dependent applications.

\subsection{The Green's function concept}\label{sec_green}
We return to Equation \cref{2.8b}, where the function $F(\vv,i;\vr'\rightarrow \vr))$ was defined.
For a finite domain, see \cref{eq1.5b}, and introducing again the shortcut $l$ for $(\vr'+l\vOmega,\vv,i)$
this becomes:
\begin{equation}
F(l) = \left\{
\begin{array}{lr}
\exp \left [   - \int_{0}^{l} ds \Sigma_t ( s
 )  \right ],                  &   l<l_{max}  \\
0                      &   l\ge l_{max}
\end{array}
\right.
\end{equation}

This function $F$ is the Green's function of the left hand side (the ``convective part") of the transport equation.
\bigskip

to be written


\newpage
\section{Monte Carlo solution of equation 1} \label{sec1.2}

A statistical solution to equation \cref{eq1b} is straight forward,
because it is formulated in probabilistic terms as follows.

A discrete Markoff chain is defined using $Q$ as an initial
distribution and $L = T  \cdot  C$ (order of $C$ and $T$ reversed
compared to $K$, $K = C \cdot T$) as a transition probability. Histories $\omega^n$
from this stochastic process are generated according to $\omega^n
=   (x_0, x_1, x_2, ... , x_n)$, (where $x_j   =   x_a $ for all $
j   \ge   n$ and $ x_i   \ne   x_a $ for all $ i   <   n$), with
$x_n$ being the first state after transition into the absorbing
state $x_a$. $ x_0$ denotes the initial state distributed as
described by $Q$. Thus, the length n of the chain $ \omega^n$ is
itself a random variable. A random sampling procedure to generate
such chains is carried out in Monte Carlo codes by converting
machine generated (pseudo-) random numbers $ \xi_{i_1} , \xi_{i_2}
, ...$ into random numbers with the distributions $Q$, $T$ and
$C$. Having computed N  chains $ \omega_i$,
$i=1,2,...,N$, the responses $R$
\cref{eq1.7} with respect to detector function $g$ are estimated as the arithmetic mean
of functions (``statistics'', or  ``estimators'') $X_g(\omega )$, i.e.
\begin{equation} \label{eq1.8}
R  \approx   \tilde{R}   =    \frac{1}{N} \sum_{i=1}^{N}   X_g(\omega_i) .
\end{equation}
The proper choice of the number of histories $N$ depends on the variance $\sigma^2(X_g)$ of the estimator $X_g$ and is highly problem specific.
In EIRENE it can range from $N=2$ for conditional expectation estimators and point sources in phase space, up to
values of several millions for $N$.
\subsection{Unbiased estimators}\label{sec1.2.a}
One possible choice for $X(\omega )$ is the so called ``collision
estimator" $X^c$, which, for a chain no. $i$,  $\omega_i^n$ of length $n$ reads:
\begin{equation} \label{eq1.9}
X^c_g ( \omega^n_i)   =   \sum_{l=1}^{n} g_c (x_l )
 \cdot   \prod_{j=1}^{l-1}
 \frac{c(x_j )}{(1 - p_a(x_j))}   \ \ ,
\end{equation}
and the summation index $l$ runs from the first point of collision $x_1$ until the last (absorption) at $x_n$.
The product in the sum accounts cumulatively for branching, and absorption rates during along the Markoff chain $\omega_i^n$.

This estimator is, for example, used in the DEGAS code
\cite{kn:Heifetz}. It can be shown \cite{kn:Spanier69a} in a tedious but mathematically strict way that the statistical
expectation E($X^c$) produces:
\begin{equation} \label{eq1.10}
R   =   E(X^c )   =   \int  d( \omega )   X^c_g ( \omega ) h( \omega
)
\end{equation}
with $h(\omega)$ being the probability density for finding a chain
$\omega$  from the Markoff process defined above, i.e. $X^c$ is an
unbiased estimator for response $R$.
\\
Other estimators (``track-length type estimators") are employed
frequently. These estimators are unbiased as well but have higher
moments different from those of $ X^c$. Instead of evaluating the
detector function $ g_c(x)$ at the points of collisions, $ x_l   ,
$ (as $ X^c$ does), they involve line integrals of $ g_t(x)$ along
the trajectories, e.g.,
\begin{equation} \label{eq1.11}
X^t_g ( \omega^{n}_i )   =    \sum_{l=0}^{n-1}   \left \{
\int_{x_l}^{x_{l+1}} ds \  g_t (s) \right \}
  \cdot   \prod_{j=1}^{l-1}
 \frac{c(x_j )}{(1 - p_a(x_j))}   ,
\end{equation}
with $ R   =   E(X^t )   =   E(X^c)$.
In literature mostly the detector function $g_t$ is taken to be constant within a grid cell.
The line integral in Eq. \cref{eq1.11} then collapses to the product $g_t s_l$, with $s_l$ being the distance
(``track-length'') travelled in cell during step $l$.
The mild generalisation of the track-length estimators
to detector functions, which can vary along the track even within a cell, is derived in
\cite{kn:Reiter84a} and applied in EIRENE as early as 1984 in \cite{kn:Reiter84b}.

It has been shown, e.g., in
\cite{kn:Spanier69a}, that the collision estimator, derived not
for the pre- but for the post-collision density integral equation
results in a track-length type ``conditional expectation
estimator" $ X^e$, which, together with $X^c$ and the
``traditional" track-length estimator $ X^t$, may be used as one
further option in the EIRENE code.
\\ This estimator $ X^e$ is obtained from $ X^t$ by
extending the line integration, which is restricted to the path
from $ x_l$ to $ x_{l+1}$ in formula \cref{eq1.11}, to the line
segment from $ x_l$ to $ x_{end}$. Here $ x_{end}$ is the nearest
point on a boundary along the test flight originating in $ x_l$.
I.e., the line integration may be extended into a region beyond
the next point of collision, into which the generated history
would not necessarily reach. This ``conditional expectation
estimator" reads:
\begin{equation} \label{eq1.12}
 X^e_g ( \omega^n_i )   =
  \sum_{l=0}^{n-1}  \left \{ \int_{x_l}^{x_{end}}
ds \  g_t (s)    \cdot   \exp \left (
 - \int_{0}^{s} ds '   \Sigma_t (
  s ' ) \right ) \right \}
   \cdot   \prod_{j=1}^{l-1}
 \frac{c(x_j )}{(1 - p_a(x_j))}   ,
\end{equation}
If $ x_{end}$ is taken to be the nearest point on each mesh cell
boundary, then the estimator \cref{eq1.12} reduces, as a special
case, to the method employed by the NIMBUS code \cite{kn:Cupini}.
However, the estimator \cref{eq1.12} is more general, as the
integration may be extended over arbitrarily many cells. The
length of the integration path is controlled in EIRENE by an input
flag WMINC, see \cref{sec2.10}.


\subsection{Scaling of tallies}\label{sec1.2.0}
With increasing $N$, the number of Monte Carlo histories,
the unbiased estimators $X$ given in previous \cref{sec1.2.a} provide arbitrarily precise approximations for the responses $R_g=\langle\Psi|g\rangle$,
for detector function $g$ and dependent variable $\Psi$ \cref{eq1d}. The precise meaning and physical dimensions of $g$ and $\Psi$
depend upon the problem at hand, e.g. also on symmetry (ignorable spatial coordinates, or time).

Because the inhomogeneous part of the transport equation \cref{eq1d} $\cal{S}$ was assumed to be normalized to one (for source
sampling),
\begin{equation}\label{scale}
{\cal{S}} (\vr,\vv,t)=\frac{1}{s} \ S(\vr,\vv,t), \quad  s =\int  d^3rd^3vdt\ S(\vr,\vv,t)
\end{equation}
this normalization factor $s$ has to be multiplied to estimators \cref{eq1.8} to turn responses $R_g$ \cref{eq1.7} into
estimates with correct dimensionality and in absolute units.

Mostly the responses of interest are \emph{intensive quantities}, such as density, pressure, collision density (number
of collision per unit time and volume) rather than
the \emph{extensive quantities} obtained by the Monte Carlo phase space integration in \cref{eq1.7}, total energy, momentum or
number of particles.  The volume, e.g. of a grid cell, is also an extensive quantity. Thus by dividing the estimate by the appropriate cell-volume $V_m$ (in space-time), see \cref{eq2.10}, finally the
profiles of cell averaged intensive quantities (e.g. density profiles) are obtained.
The final unbiased estimate, in absolute units, for any of the unbiased estimators $X_g$ for detector function $g$ discussed above, then is:
\begin{equation} \label{eq1.8b}
R_g  =\langle\Psi|g\rangle \approx   \tilde{R}_g(N)   =    \frac{s}{ V_m\times N} \sum_{i=1}^{N}   X_g(
\omega_i) , \quad N \gg 1
\end{equation}
with $N$ being the number of Monte Carlo histories.
Correct interpretation of the results of an EIRENE run hence requires knowledge of the
ratio $s/V_m$.
The total source strength scaling factor $s$, i.e. the integral of the
inhomogeneous part $S$ of the governing Fredholm integral equation, is input to EIRENE,
variable ``FLUX" in input block 7 (see \cref{sec2.7}). The space-time cell volumes $V_m$ are either automatically calculated by EIRENE
in subroutine ``VOLUME", for cells belonging to standard grid (input block 2, \cref{sec2.2}), or are provided from external considerations
by a number of options, e.g. input blocks 3b, 5, or 8.
In this latter case great care is needed, however, that the cell volumes are proportional to the volumes
\emph{as seen by the test flights} for otherwise not only unscaled profiles, but even wrong (biased) profile shapes will be obtained.

\subsubsection{Scaling in problems with ignorable coordinates}\label{sec1.2.0a}
As seen in the previous paragraph, the absolute values of estimates are scaled
with the ratio $s/V$ of source strength to cell volume. Depending on ignorable coordinates in any particular problem,
or on whether stationary (i.e.: time is an ignorable coordinate) or time dependent transport problems are considered ,
the dependent variable $\Psi$
in the governing integral equation \cref{eq1d},
the source strength $s$ and ``cell volume" $V$ may have different interpretations:
\paragraph{stationary (time independent) problems}
\paragraph{one ignorable spatial coordinate}
In stationary problems with one ignorable coordinate, say, the z-coordinate,
the source strength $s$ (input flag ``FLUX") is the flux, particles per
unit length $dz$ in direction of $z$. Likewise, the volume $V$ is per unit length in the ignorable direction,
i.e. if $dz=1$ then $V$
is the cell area in the two remaining coordinates $x,y$.
The unit length $dz$ of an ignorable coordinate (here: $z$-coordinate) used in a particular run
is determined by the input flags in the corresponding input block for standard grid options,
i.e. in input block 2A for the x-coordinate, input block 2B for the y-coordinate and input block
2C for the z-coordinate, see \cref{sec2.2}. Note that the numerical value of source strength
``FLUX" (block 7) corresponds to the choice of e.g. $dz$. If $dz=1$ \si{\centi\metre}, then ``FLUX" is the number of particles
per unit time and per \si{\centi\metre} in z-direction. If $dz=1$ m, then the same value of variable ``FLUX"
would correspond to the 100 times smaller source strength of ``FLUX" particles per unit time and per meter.
All resulting volume averaged output tallies would have a value 100 times smaller. A particle density
might then also be interpreted as surface density (particle per unit area).



\subsection{Statistical errors, Efficiency (\acs{FOM})} \label{sec1.2.1}
The efficiency for Monte Carlo Codes is the
inverse of the \ac{FOM} of the calculation, defined
as
\begin{equation} \label{eq1.13}
\texttt{FOM} = statistical \ variance    \cdot   computing \ cost.
\end{equation}
Note that \texttt{FOM} should be approximately independent of the running
time, because the number of histories generated is (excluding
overhead) proportional to the CPU costs and inversely proportional
to the statistical variance $\sigma^2$.

It is one of the major advantages of Monte Carlo methods over other numerical
schemes that the error estimates, empirical variances $\tilde\sigma^2$, are directly provided by the method itself,
not requiring any further considerations.
The options to activate evaluation of statistical variance in an EIRENE run for any computed quantity
(tally) are described in input block 9 (\cref{sec2.9}). EIRENE provides numerical and graphical
output for the ``empirical relative standard deviation", in \%.
\begin{equation}
\tilde\sigma_{g,rel}(N) = \tilde\sigma_g(N)/\tilde R_g(N) \times 100
\end{equation}
and $\tilde R_g^N$ is the Monte Carlo estimate for tally $R_g$  based on $N$
Monte Carlo histories \cref{eq1.8b}, properly scaled to source strength $s$ \cref{scale}. Here $R_g$ may be any volume or surface averaged tally
for detector function $g$, either track-length, collision, or snapshot estimated.
The  variance $\sigma^2_1$ per history is obtained as the (unbiased) estimate (``empirical variance") as:
\begin{equation}\label{MCsig1}
\sigma^2_1 \approx \tilde\sigma^2_1(N) = \frac{1}{N-1}\sum_{i=1}^N (X_i-\overline{X})^2=\frac 1{N-1}\left[\sum_{i=1}^N X_i^2 - \frac 1N \left(\sum_{i=1}^N X_i\right)^2\right]
\end{equation}
where $X_i= X_g(\omega_i)$ is the contribution of Monte Carlo history $\omega_i$ to
estimator $X_g$ for tally $R_g$, and
$\overline{X}$ is the arithmetic mean $\overline{X}=1/N\sum X_i$ over all histories. $N$ is the number of (statistically independent)
Monte Carlo histories. The subscript $g$ is omitted here and from now on. Note that the expression on the right hand side of \cref{MCsig1} can be evaluated after each completed particle history,
i.e. without storing all individual contributions $X_i$ first.

The variance per history is turned into the final estimate for the Monte Carlo variance by the ``law of large numbers" (and the ``central limit theorem"
of probability theory):
\begin{equation}\label{MCsig}
\tilde\sigma^2_{MC}(N)=\frac 1N \tilde \sigma_1^2(N)
\end{equation}
For large $N$ the variance per history (= variance of a single observation) $\tilde\sigma_1^2$ converges to a  constant (namely to $\sigma^2_1$),
and the final Monte Carlo estimate of variance $\tilde \sigma^2_{MC}$ \cref{MCsig}
has the expected probabilistic $1/N$ scaling,
i.e. the empirical standard deviation, which is the Monte Carlo error estimate,  scales as: $\tilde\sigma_{MC} \sim 1/\sqrt{N}$.
\paragraph{Multiple cell crossings}

One should note that except in simple 1D cases in general one single Monte Carlo history $\omega_i$ can contribute more than
once to the estimate $\tilde R_g$ for a particular cell of the grid or surface segment. E.g. test particles can cross one cell $m$ [see \cref{eq2.10b}] more than once,
with each cell crossing ``$j$" leaving a score (contribution) $X_{i,j}$ to estimator $X_g(\omega_i)$. Then:
\begin{equation}\label{mccross}
X_i=\sum_{j=1}^{J_i} X_{i,j} \quad ,
\end{equation}
with $J_i$ being the number of contributions  (e.g. cell crossings) to the estimator, from particle history no. $i$.
The summation in \cref{eq1.8} and hence also the second sum on the right hand side of  \cref{MCsig1} can still be accumulated
``on the fly" while generating the Monte Carlo histories (i.e. at no extra CPU cost).
However, because, of course  $X_i^2 \neq  \sum X_{i,j}^2$, the first sum on the right hand side of \cref{MCsig1}
cannot be evaluated ``on the fly" (per event). Instead here EIRENE first accumulates the sum \cref{mccross} $X_i$ for individual test flight ``$i$", and updates the
sum $\sum X_i^2$ only once after each completed history (in subroutine. STATIS, entry STATS1). I.e., variance estimates are made after completion of each history, inside the particle loop, see \cref{flow5}.
The full evaluation and scaling of
expression \cref{MCsig1,MCsig} is then carried out after completion of all histories (per stratum), in entry STATS2 of subroutine STATIS.
\paragraph{Linear functions of tallies}\label{lin-funct}
Often algebraic functions of tallies (sums, products, ratios,....) are formed after a Monte Carlo run to produce further (post processed) output quantities
(see input blocks 10C and 10E in \cref{sec2.10}).
For those tallies standard deviations are not available in general, due to statistical correlation between individual tallies obtained from the same Monte Carlo run.

However, for linear functions of tallies (e.g. sums, differences), this is possible even without resorting to the covariance estimators (block 9, \cref{sec2.9}) of EIRENE.
Usually this was done in ``problem specific" modules, e.g. in STATS1\_COP for source terms arising in plasma fluid codes (see \cref{sec2.14}).
In 2012 a new routine (UPFCOP) was added, and
is called inside the particle loop, scored per particle (not per event) and prior to STATS1.

Let a linear function $R_{lin}$  of default tallies $R_k$ be given by
\begin{equation}
R_{lin} = \sum_k a_k R_k,
\end{equation}
($a_k$ are some constant scalar factors)
then after each completed history $\omega_i$ in routine UPFCOP the sum
of contributions
\begin{equation}\label{interchange}
X_l(\omega_i)= \sum_k a_k X_k(\omega_i) = \sum_k a_k \sum_j X_{k,j}(\omega_i) = \sum_j \sum_k a_k X_{k,j}(\omega_i)
\end{equation}
is formed, to finally build the linear function tally $R_l(N)=1/N\sum_i X_l(\omega_i)$ by averaging over all histories, same as for default tallies.
This is possible because in case of linear combinations of tallies the summation over tallies $k$ and event-contributions $j$ of a single
history can be interchanged \cref{interchange}.


Finally the variances for $R_l$ are evaluated in STATS1, reusing the procedure used for default tallies, and the
former problem specific routine STATS1\_COP for variances from tallies needed for code coupling has hence been made redundant, as indicated in red color in \cref{flow5}.
Currently linear combination tallies $R_l$ are stored
on array (tally) COPV. For them all printout and graphical output options are available.
One example of such linear combination tallies is discussed in \cref{sec2.14} for interfacing EIRENE with
\ac{CFD} codes, as usually the (kinetic) reaction source terms provided by EIRENE to plasma fluid codes are such linear functions
of default tallies.

\paragraph{Statistical independence of MC histories}

When the Monte Carlo histories from one EIRENE run are not strictly statistically independent,
as e.g. in case of stratified source sampling (see paragraph below), then a modified
formula for the statistical error estimates has to be used.

\subsubsection{Sampling, Non-analog sampling}
\label{sec1.2.1a}

One often encounters in literature the description of special,
``new" Monte Carlo techniques, that are greatly superior to so
called ``standard methods".

It is true that one can devise very intelligent methods to
optimize performance, but each optimization always only works well
for one very particular problem, it can equally well entirely
wreck the performance of an only  slightly different case.

For a general purpose Monte Carlo solver for transfer problems, as
EIRENE, therefore, no general advise can be given, although the
performance often could greatly be improved by adapting the method
to a particular problem. In particular, non of these ``intelligent
methods" have been (and will ever be) hard wired into EIRENE.

EIRENE contains a set of such methods, referred to as
``non-analog" methods and controlled by the flags in input block
9, see \cref{sec2.9}. Activating those must be accompanied
by a very careful statistical analysis of results, not just the
run time per particle, nor even the standard deviation alone,
suffices.


We first note here that any random distribution function $f(x)$
arising during the course of generating the chains $\omega$
(``particle histories") can be replaced by another ``non-analog"
function, $ g(x)$, if a weighting factor,
\begin{equation} \label{eq1.13a}
 w(y)   =   f(y) / g(y)
\end{equation}
is included in the formula for the estimator $X(\omega )$, y being
an actual random number generated from $ g $. This choice
can in some cases increase the efficiency of the algorithm.
\\
The only restrictions on the choice of non-analog sampling functions $g(x)$
(besides practical ones) are:

\begin{equation} \label{eq1.14a}
\mbox{if} \quad  g(x) = 0, \qquad \mbox{then}  \quad f(x)
= 0;
\end{equation}

and conditions to ensure that the non-analog process remains
sub-critical as well.

The condition \cref{eq1.14a} is checked in an EIRENE run whenever
non-analog distributions are applied, and, if violated, an error
exit with the message: ``violation of Radon-Nikodym condition"
occurs.

Note: a violation of condition \cref{eq1.14a} can not be detected
otherwise, e.g., by monitoring overflow in the weighting factor $w$.
Of course, values y resulting in a zero in the denominator of $w$
\cref{eq1.13a} have sampling probability zero according
to distribution $g(x)$, and, hence, are never sampled.
Still the results would be biased.

Note: when inspecting an EIRENE geometry plot with particle
trajectories (or even a movie: NLMOVIE=TRUE, input block 1) in
order to get an intuitive feeling for the particular transfer
process considered, then all non-analog sampling options must
first be turned off, (NLANA flag in input block 1), for otherwise
the pictures may be grossly misleading.

\medskip
We consider the procedures for random sampling from univariate
distributions as known, and refer to the many textbooks on that,
in particular to the ``random sampling library"
\cite{kn:Sampling}. If none of the direct methods apply, then
still either the non-analog method mentioned above, or the
``rejection method" can be used.

Sampling from a multivariate distribution $f(x_1,x_2,...,x_n)$ can
always be reduced to a sequence of samplings from univariate
distributions, by noting that:
\begin{equation} \label{eq1.15a}
f(x_1,x_2,...,x_n)=f_1(x_1) \cdot f_2(x_2|x_1) \cdot
f_3(x_3|x_1,x_2) \cdot ...
\end{equation}
Here, $f_1$ is the marginal distribution obtained from $f$ by
integrating over all but the first independent variables. It is a
univariate distribution of $x_1$.

$f_2$ is a conditional marginal distribution obtained from $f$
firstly by integrating $f$ over all but the first 2 variables
$x_1, x_2$, and secondly then taking the conditional distribution,
conditional on $x_1$ (i.e., the univariate distribution of $x_2$
for given values of $x_1$).

Likewise, $f_3$ is a conditional marginal distribution, n-3 fold
integration of $f$, and then distribution conditional on $x_1, x_2$,
and so on.

Random sampling from $f(x_1,x_2,...,x_3)$ then proceeds by
sampling first $x_1$ from the (univariate) first factor in
\cref{eq1.15a}, then $x_2$ by sampling from the
(univariate) second factor, and so on.

This transformation of a multivariate (k-) distribution to a uniform distribution
on a k-hy\-per\-cube is sometimes referred to as ``Rosenblatt transformation" (M. Rosenblatt, Ann. Math. Statistics 23: 470-472, (1952)).

One particular example of this scheme is the ``TRIM database
surface reflection model" (see below, \cref{sec1.3.2}).
There, essentially, tables of the conditional marginal
distributions mentioned here are pre-computed with the TRIM code.
For convenient random sampling, these tables are stored for the
inverted cumulative distribution, i.e., as conditional quantile
functions.

\subsubsection{Stratified Source Sampling}
\label{sec1.2.1b} The EIRENE code resorts to a ``stratified
sampling" technique. This technique is one the few Monte Carlo variance reducing techniques
which are quite straight forward, easy to implement and, most importantly,
behave in a well predicable way. It is therefore recommended to always use ``source stratification"
described below as much as possible.
There are, however, non-trivial issues of load balancing in case of multi-core runs (i.e. Monte Carlo code parallelisation),
which are related to assignments of strata to cores, in particular if the average cpu-time per history
strongly varies between strata. In such cases the ``proportional allocation of weight" to source strata discussed below is
technically difficult to achieve together with proper load balancing, due to the stochastic termination of the random walks
in the various strata. Whereas without stratification extracting parallelism seems as trivial
as constructing many independent samples in parallel, this ``embarrassingly parallel'' feature of linear Monte carlo transport codes is lost, and
parallelisation in combination with stratified source sampling  requires some attention.




Stratified sampling means that the primary source
distribution $Q(\vr,\vv,t)$ , see Equation \cref{eq1a}, or more generally, for $S(x)$ in \cref{eq1d}, can be decomposed into a sum of
$M$ independent sources (this applies for both: the source distribution $S$ itself and the corresponding source strength $s$):
\begin{equation}\label{sss}
S = \sum_k^{M} S_k ; \quad s=\sum_k^{M} s_k
\end{equation}
and the solution is obtained
by linear superposition of the solutions for each sub-source
(``stratum") $S_k$. We use again the slightly more general notation $S$ for the inhomogeneous part of the Boltzmann equation
rather than $Q$, which is the physical source of particles (objects), i.e. a special case of $S$.

The stratified sampling technique is a well known statistical
procedure from experimental planning, see any textbook on statistics or Monte-Carlo methods.

It sometimes can significantly affect the efficiency of a
Monte-Carlo run. This can go both ways, depending upon the
particular problem and the particular stratification used.
Therefore, as with the non-analog options, no general
recommendation can be made.

The effects of source stratification, however, are usually more easy to
assess (predict) than those resulting from general non-analog methods.

In order to make connection with the textbooks, in which ``stratified sampling" is often discussed
in connection with Monte Carlo integration procedures,
the Monte Carlo solution of the transport problem is interpreted as a Monte Carlo integration,
and the ``stratified sampling" in Monte Carlo integration becomes a ``stratified source sampling" in
transport problems in this interpretation. To see this we need to note
here only the following (based on Fredholms alternative for the linear Boltzmann operator):

EIRENE provides estimates of linear functionals (``moments",
``responses")

\begin{equation}  \label{fredhalt}
\langle g,\phi_S\rangle    = \langle S,\phi^*_g\rangle
\end{equation}

As described above, $g$ is the detector function (``weighting
function") $\phi_S $ is the solution to the kinetic transfer
equation (Boltzmann equation) in which $S$ was used as (inhomogeneous)
source term. $\phi^*_g$ is the solution of the corresponding
adjoint equation, with $g$ as source term (and $S$ as ``weighting function".

The above equation means that EIRENE estimates can simply be
regarded as Monte-Carlo estimates of multi-dimensional integrals,
namely of the adjoint function $\phi^*_g$ integrated over phase
space with $S$ defining a distribution (measure) in phase space ``$\mu_S$".

It is then obvious that one can split that integral into a sum of
integrals by arbitrarily decomposing the domain of integration,
i.e., by decomposing the source $S$ into sub-sources (``strata")
$S_k$ with
\begin{equation}
R_g = \int d\vx~ \phi_S \cdot g    = \langle S,\phi^*_g\rangle  =\sum_k^{M} \langle S_k,\phi^*_g\rangle = \sum_k^{M} \int d\mu_{S_k} \phi^*_g
\end{equation}
This is the concept of ``stratified sampling", which, for Monte Carlo
particle transfer procedures in particular actually turns out to
be a ``stratified source sampling".

Further details about this in our particular context are given in
the description of input block 7, \cref{sec2.7} below. As pointed out above
the evaluation of statistical error estimates has to account for the stratification,
because with stratification  the Monte Carlo histories are not strictly independent anymore.

Grouping the $N$ statistically independent test flights (random samples) into $M$ strata, leads to Monte Carlo particle
numbers $N_k$ sampled from stratum $k$ (i.e., with birth points sampled only from source $S_k$) with:
\begin{equation}\label{Nstra}
\sum_{k=1}^{M} N_k = N   \quad , \quad  N_k > 0 ~ \quad \mbox{for}~~ ~ k=1,...M
\end{equation}
where, again, we omit the subscript $g$ for the particular tally (response function) from now on.

The choice of the $N_k$
is only restricted by the normalization condition in \cref{Nstra}.
This therefore also applies for the allocation of CPU time to individual strata for a given overall CPU time.
This additional freedom can be used in EIRENE runs to affect the statistical error for a given CPU time.
Unfortunately, as with non-analog sampling, this can go both ways. Fortunately, however, here a clear general
recommendation can be given, see below in this section: ``proportional allocation".

Lets denote by
$\tilde R_k(N_k)$ the estimate for tally $R$ as obtained only from the $N_k$ samples from stratum $k$ (but scaled with
the full source strength $s$). Likewise, we denote by
$ \tilde \sigma^2_{1k}(N_k)$ the variance of estimate $\tilde R_k$ per history, i.e. for a single sample,
obtained by only using the $N_k$ samples from stratum $k$ (again implying the scaling with total source strength $s$).
 Then clearly
\begin{equation}
\tilde R^*(N) =\sum_k  \frac {s_k}{s} \tilde R_k(N_k)
\end{equation}
is the (only) unbiased weighted estimate of $R$ obtained from summing over all strata, i.e. from all $N$ samples.
The asterisk denotes the estimate $\tilde R$ of $R$ obtained with stratification, $N=N_1+\dots + N_M$, and $\tilde R(N)$ is the estimate
of the same quantity from a simple sample (no stratification) of same size $N$:

\begin{equation}
R = E\{\tilde R^*(N)\} = E\{\tilde R(N)\}
\end{equation}
where $E\{X\}$ denotes the ``expectation value" of random variable $X$.
Furthermore:
\begin{equation}\label{varstra}
\tilde \sigma^2_{k}(N_k)  = \frac{1}{N_k}\tilde \sigma^2_{1k}(N_k); \quad \tilde \sigma^{2^*}(N)  = \sum_k \frac{s_k^2}{s^2}\frac{1}{N_k}\tilde \sigma^2_{1k}(N_k)
\end{equation}
where again the asterisk indicates the value obtained with stratification.
The right equation follows from the sum rule for variances and because
the estimates for the strata are mutually statistically independent.
It is clear that the variance with stratification $\tilde \sigma^{2^*}(N)$ can certainly be different from the variance $\tilde \sigma^2(N)$ obtained  from a simple random sample (without stratification) of the same size $N$.
To see this we write the trivial identity:
\begin{equation}\label{varsimp}
\tilde \sigma^2(N) = \sum_{k=1}^{M}\frac{s_k}{s} \tilde\sigma^2(N)
\end{equation}
For example, in case of ``proportional allocation" of CPU-time to strata according to their source strength $s_k$ , i.e.,  $N_k \propto s_k$ :
\begin{equation}\label{propalloc}
N_k=\frac{s_k}{s} N, \quad k=1,2,\dots,M
\end{equation}
one finds from \cref{varstra} for the variance of the stratified sample:
\begin{equation}
 \tilde \sigma^{2^*}(N)  = \sum \frac{s_k}{s}\tilde \sigma^2_{k}(N_k)
\end{equation}
which can be substantially smaller than the variance from the simple sample \cref{varsimp} of size $N$, if the
subdivision into strata is made such that the variability within strata $\sigma_k^2$ is less then the variability
in the entire population $\sigma^2$.
Furthermore, one can show that with proportional allocation \cref{propalloc} the variance from the stratified sample is
always less than or equal to the variance from the simple sample of same size $N$, no matter how the subdivision into strata was
done:
\begin{equation}
\mbox{If} \quad N_k=\frac{s_k}{s} N; ~~ k=1,2,\dots,M \quad \mbox{then} \quad  \tilde \sigma^{2^*}(N) \leq \tilde \sigma^{2}(N)
\end{equation}
i.e., stratified source sampling with proportional allocation is an ``inherently safe working procedure" for variance reduction.

This important inequality can be derived after some lengthy algebra, leading finally to
\begin{equation}
\tilde \sigma^{2}(N)=\tilde \sigma^{2^*}(N)+ \frac{1}{N}\sum_k[ R_k(N_k)- R(N)]^2
\end{equation}
Clearly, the second term on the right side of this equation, which is the variance between strata, is always positive (or zero)
and this fraction of the variance is eliminated from the Monte Carlo estimate by stratified sampling.

This fact is the reason for the EIRENE option of ``proportional allocation of CPU time" described in input block 7,
\cref{sec2.7}.

In principle one can even optimize the CPU performance (i.e. minimize the variance per CPU time, for any given stratification $S_k$) by choosing the
allocation $N_k$ to strata in an even more intelligent way (involving a-priori knowledge of the variability within strata $\sigma^2_k$,
e.g. estimated from a previous iteration.)  But these options are likely to be quite sensitive to the choice of the particular tally (response $R_g$).
In a typical EIRENE application many such responses for a large variety of detector functions $g$ are estimated simultaneously in one single run.
Therefore this ``optimal allocation" schemes are not currently implemented.

\subsubsection{Testing the stratification}
The proper implementation of the stratification, for all kinds of output tallies (such as volume averaged tallies,
surface averaged tallies, spectral tallies) can be verified by including two:  $M=2$ (or more: $M$) identical strata,
even with identical number of test-histories $N_k$ and identical random number seeds for those.
The standard deviations of the sum over these identical strata
must then be smaller by a factor $\sqrt{2}$  (or: a factor $\sqrt{M}$, respectively) \emph{exactly}, rather than
only approximately, as it would be the case without stratification by simply enhancing the number of test-histories
for a stratum by a factor two (a factor $M$, respectively).

\subsection{Source sampling} \label{sec1.2.2}
Source sampling in EIRENE is carried out in routine LOCATE.F. This routine converts uniform random numbers into
random numbers with distribution $S_k$, the k-stratum part of the inhomogeneous term $S$.
Again sequential sampling \cref{eq1.15a} from conditional distributions is applied.

In LOCATE.F the following sequence is used (stratum no. $istra=k$ is set in calling program):
\begin{equation}
S = S_1(\vx) \times S_2(i|\vx) \times S_3(\vv|\vx,i)
\end{equation}
with $S_1(\vx)= \sum_i\int d^3v S(\vx,i,\vv)$ being the cumulative distribution of the spatial coordinate $\vx$ of the ``birth point"
of the trajectory, obtained by integrating the source term  $S$ over all velocities and summing over all species.
$S_2(i)$ is the distribution of birth species index, obtained by integrating over all velocities and conditional on the already
sampled spatial birth point.
$S_3(\vv)$ finally is the velocity distribution at the birth-point of the trajectory, conditional on the already sampled
species index $i$ and birth point location $\vx$.

This procedure is employed for all types of sources, namely for re-sampling from a stored ``census array" (e.g. initial condition
in time dependent mode), for point sources, line sources, surface sources and volume sources.

If the sampled species $i$ belongs to the background species ($i=i_b$) rather than to the test particle species community, then, after completion
of sampling from $S$ an
appropriate ``collision" is carried out (e.g. with a wall surface in case of surface sources), until a test particle species $i_t$ results.
I.e., formally a sampling from  kernel $C(\vx,\vv,i_b \rightarrow i_t,...)$ is carried out at the place of birth $\vx$ to produce initial test particle coordinates.

Stratification of sources has been described in \cref{sec1.2.1b}.

\subsubsection{time-census sources}
\subsubsection{Point sources}
\subsubsection{Line sources}
\subsubsection{Surface-sources}
and see special section on ``plasma surface recycling source
sampling", \cref{sec1.4}.
\subsubsection{Volume sources}
special case: volume (radiative or threebody) recombination of plasma ions and electrons to neutrals.
Also used for: spontaneous photon emission source in radiative transfer simulations.

\subsection{Sampling a free flight from the transport kernel T} \label{sec1.2.3}
The conditional sampling distribution $T(\vr'\rightarrow \vr)$ for the next point of collision (event), given a test flight
is at position $\vr'$, traveling with velocity $\vv'$, is given by \cref{eq1.5}.

A more intuitive argument for this distribution of free flight path lengths follows from the
simple attenuation law of a mono-energetic beam incident on a thin slab of material, thickness $dl$,
target material density $n_t$. Clearly the fraction of beam particles undergoing collisions in this slab
is $\sigma n_t dl$, with ``cross section'' $\sigma$ (stationary target). This  results in the attenuation law
\begin{equation}
dn = -n n_t  \sigma dl \quad \quad \rightarrow \quad n = n_0 \exp(-n_t \sigma l)
\end{equation}
where $n$ is the number of beam particles, $n_0$ the number of beam particles at $l=0$.
It is then naturally supposed that
\begin{equation}
p(l)dl = [\exp(-n_t\sigma l)]n_t\sigma (dl)
\end{equation}
is the probability for a first collision between $l$ and $l+dl$,
and hence
\begin{equation}
P(l) = \int_0^l [\exp(-n_t\sigma l)]n_t\sigma (dl) = 1 - \exp(-n_t \sigma l)
\end{equation}
is the corresponding probability distribution function for a first collision at distance $\le l$,
i.e. the same results as derived
above formally with the Green's function argument.

Returning to \cref{eq1.5}
we first  again omit velocity and species index in the argument lists to keep the notation simple, but we note that certainly
the mean free path $\lambda$ (or macroscopic cross section $\Sigma_t$) usually depend explicitly on the test particle energy and species.
Let us introduce the distance of the flight $l$ in direction $\vOmega=\vv/|\vv|$, hence
$\vr=\vr'+l\vOmega$, then the sampling distribution for $l$ is simply
\begin{equation}
T(l)=\Sigma_t(\vr'+l\vOmega)\exp\left[-\int_0^l ds \Sigma_t(\vr'+s\vOmega)\right], \quad \Sigma_t=1/\lambda_t
\end{equation}
Distances $l$ are readily sampled from this exponential distribution using the
inversion method and the cumulative distribution $G_T$ of $T$
\begin{equation}\label{cummul}
\xi=G_T(l)=1-F(l)=1-\exp\left[-\int_0^l ds \Sigma_t(\vr'+s\vOmega)\right]
\end{equation}
with $\xi$ (and hence also $1-\xi$) being a uniformly distributed  random variable on [0,1].
Therefore, cumulating $\left[-\int_0^l ds \Sigma_t(\vr'+s\vOmega)\right]$ along a flight and comparing it with
$\ln(\xi)$ with a uniform random number $\xi$ provides (by inversion) a random sample of the free flight distance $l$
with the proper distribution.

Sampling the distance $l$ of a free flight, for test particle species ``$i$" with velocity $\vv$, from the transport
kernel $T$ requires knowledge of the (``monochromatic") mean free
path $\lambda = \lambda(\vr,|\vv|)$ along the
trajectory, and, hence, of the ``macroscopic cross section"
$\Sigma = \Sigma(\vr,|\vv|)$, which also may
vary along the trajectory.

The ``macroscopic cross sections" $\Sigma$ (by abuse of language,
dimension: 1/length) with or without subscript $k$ (the label for
the type of the collision process) or subscript $t$ (for ``total", the sum
over $k$), are defined as:
\begin{equation}
\Sigma = \frac{1}{\lambda} = \frac{\nu}{|\vv|}, \qquad
\lambda = \frac{|\vv|}{\nu}
\end{equation}
with the mean free path $\lambda$ and the collision frequency
$\nu$:
\begin{equation}
\nu(\vr,\vv,i) = n_b \langle\sigma \cdot
v_{rel}\rangle_b, \qquad v_{rel}= |\vv-\vv_b|
\end{equation}
Here $n_b$ and $\vv_b$ are the target particle density
and velocity, respectively, and the brackets
$\langle...\rangle_b$ denote averaging with the velocity
distribution $f_b(\vr,\vv_b)$ of the
background medium species $b$.

$\sigma$ is the total scattering cross section, $\sigma =
\sigma(|\vv_{rel}|)$.

Hence, e.g., for all isotropic target distributions $f_b$ one has:
\begin{displaymath}
\nu = \nu(\vr,E,\not\! \vOmega,i) \qquad
\mbox{and} \qquad \Sigma = \Sigma(\vr,E,\not\!\vOmega,i).
\end{displaymath}

Furthermore, let $f_b$ be a drifting Maxwellian, with temperature
$T_b$ and drift $\vV_b$: $f_b =
f_b(T_b,\vV_b)$.

Then:

$\nu = \nu(T_b,|\vv-\vV_b|,i)$

Finally, if the thermal speed $v_b^{therm}$ in this Maxwellian:
$v_b^{therm} = \sqrt{T_b/m_b}$ is very large compared to
typical values of $|\vv-\vV_b|$ (i.e.: nearly
incompressible flow conditions), then:


$\nu \approx \nu(T_b,0,i) = n_b \langle\sigma \cdot v\rangle_b$

and here $\langle\sigma \cdot v\rangle_b$ is the usual
``Maxwellian rate coefficient'' for the collision process, taken at
temperature $T_b$ of the background particles (species $b$) and at
zero velocity of the test particles (species $i$). For example, in
case of neutral atom or molecule collisions with electrons this
approximation is usually valid.
From the point of view of electrons or ions, colliding with a cold background of neutrals or ions ($T_b\approx 0$,
we have the opposite case: $|\vv|\gg v_b^{therm}$ and then the rate coefficient reduced to
the simple product $\sigma(v) v$ with $v$ being the test particle velocity.

\subsubsection{Alternative: fixed time step (or constant path length increment)}
If the particle orbit between two events (the Green's function) is not given explicitly,
e.g., not as simple straight line, then the orbit has to be integrated numerically, using small
time steps $\Delta t$ or path length increments $\Delta \vx$.
The algorithm to determine the next point of collision then is different:

After a free flight  time  $\Delta t$ the probability of at least one collision during this period is the function
$G_T(l_0)$ defined above, with $|\Delta \vx| = l_0 = \int_0^{\Delta t} \vv(t) dt \approx v \Delta t$ the fixed path length for this step.
This follows from
the definition of $G_T$ as cumulative distribution for the free flight length \cref{cummul}.
Clearly $0\le G_T(l) \le 1$ for all positive $l$, by definition.
Lets now
assume constant mean free path $\lambda$, and constant velocity $\vv$ during time step $\Delta t$, and $l_0=|\vv| \Delta t$, to simplify notation.

The procedure to find the free flight length (or time) is now as follows:

Rather than sampling $l$ by inversion, as in Equation \cref{cummul},
now a uniformly distributed random number $\xi$ on [0,1] is compared with the probability $G_T(l_0) = 1-\exp[-l_0/\lambda]$.

If $\xi > G_T(l_0)$, no collision has happened during this step, and the next step is carried out.

If $\xi \le G_T(l_0)$, then at least one collision has happened prior to completion of the step.

Strictly, in this case the exact position of the first collision during this step has to be identified,
by inversion, the time step has to be shortened and the particle has to be moved to the sampled point of collision in the
interval $\Delta \vx$. However, various approximations are now usually made, to simplify coding.
\begin{itemize}
\item{} Ignore the fact that more than one collisions could have happened during $\Delta t$.
\item{} replace $G_T(l_0)$ by $l_0/\lambda ~ (= \Delta t v \Sigma_t(\vv)$, see above) when comparing with uniform random number $\xi$
\end{itemize}
The second assumption results from Taylor expansion for sufficiently small $\Delta t$ such that $l_0/\lambda \ll 1$.

The first assumption is more difficult to justify generally:

If the post collision test particle species is the same  and its new velocity is not too different from that prior to the collision,
then the probability for $n$ collisions in the same time $\Delta t$ is roughly $P_n=G_T(l_0)^n$
and the number of missed collisions in $\Delta t$ is
\begin{equation}
\sum_{i=2}^\infty P_i = \frac{G_T(l_0)^2}{1-G_T(l_0)}
\end{equation}
For sufficiently small $\Delta t$ (and hence $l_0$) this error can be made small.

If, however, collisions involve change in species, and/or significant changes in velocity, then
the smallness of $l_0/\lambda$ must be guaranteed not only for the particle which is currently being considered,
but for all particle species in the system and for all possible energies (of secondaries).
E.g.: the excitation cross section for process $A  \rightarrow A^*$ may be small, hence $l_0/\lambda \ll 1$ for species
$A$, but the de-excitation (quenching) of state $A^*$ may be very rapid. Then the quenching time of species $A^*$
determines the permitted time-step for species $A$ in this approximate procedure. Other critical examples may be systems
composed of both rather stable chemical species and short living radicals, e.g hydrocarbons.

The EIRENE code therefore only uses the correct full sampling procedure \cref{cummul}, even when particle orbits (e.g. of ions) are
integrated numerically with time-steps.



\subsection{Sampling from the collision kernel C} \label{sec1.2.4}

The general rules for sampling from multivariate distributions
outlined in \cref{sec1.2.1}, Equation \cref{eq1.15a} also
apply here. In the EIRENE code sampling from the collision kernel is carried out in subr. COLLIDE, which is called when a collision is
known to have occurred at a pre-collision particle state vector [position, velocity, species (and, optional: time)] $( \vr ' , \vv ' , i ')$
by a sequence of random decisions regarding the type of collision process $k$, the species index of the post collision test particle $i$,
and the new post collision velocity $\vv$.  Collisions are instantaneous in time.
Branching and/or absorption is accounted for by statistical weights.

It is strictly not necessary, but usually most convenient and intuitively
most clear, to decompose the kernel $C$ according to the physical
nature of the various collision processes taken into account.

The collision kernel $C$ was written as (see \cref{eq1.3})
\begin{displaymath}
C( \vr ' , \vv ' , i '   \rightarrow
\vv , i) =   \sum_{k} {p_k  C_k ( \vr ' ;
              \vv ' , i '   \rightarrow \vv, i)} \
, \ p_k   =    \frac{\Sigma_k}{\Sigma_t}
\end{displaymath}
with summation over the index $k$ for the different types of
collision processes under consideration and $p_k$ defined as the
probability for a collision to be of type $k$. The normalizing
factor
\begin{displaymath}
c_k (x ')   =
  \sum_{i}   \int d \vv \ C_k (
\vr ' , \vv ' , i ' \rightarrow
\vv , i ) \   , \  {\cal{C}}_k   =    \frac{1}{c_k}
C_k
\end{displaymath}
was the mean number of secondaries for this collision process $k$. The normalized
function $ {\cal{C}}_k$ then is a conditional probability
density for post-collision state $(\vv,i)$, given the pre-collision
coordinates $(\vr',\vv',i')$.

Sampling from $\cal{C}$ (given that a collision of some kind has
occurred) proceeds most conveniently by firstly sampling the type
$k$ of the process from the discrete distribution
\begin{displaymath}
(p_k, \quad k=1,..,K),
\end{displaymath}
and then by sampling the post-collision state of the test-particle
from ${\cal{C}}_k$. Finally the weight of the test-flight after
collision is increased (or reduced) by multiplying it with the
normalization constant $c_k$.

The decomposition of $C$ into sub-kernels $C_k$ is somewhat
arbitrary, and the criteria may also be computational aspects
rather than the different physical nature of the various processes
taken into consideration.

Sampling $(\vv,i)$ from ${\cal{C}}_k$ can be done along
the same idea: first factoring out the discrete distribution for
the species index $i$, and then by sampling $\vv$ from
${\cal{C}}_{k,i}$, the conditional distribution for
$\vv$ given the type of the process is $k$ and the
species of the post collision trajectory is $i$. Let $(n_{k,i},
i=1,...,I_k)$ be the number of post collision particles of
species $i$ from process $k$, and $n_k = \sum_{i} n_{k,i}$.
I.e., omitting pre-collision state parameters $(\vr ' , \vv ' , i ')$:
\begin{displaymath}
n_{k,i} = \int d \vv \ C_k (
\vr ' , \vv ' , i ' \rightarrow
\vv , i )
\end{displaymath}
Then, normalizing,
\begin{displaymath}
(\tilde{n}_{k,i} \quad i=1,...I_k)
 = (n_{k,i}/n_k, \quad i=1,...I_k)
\end{displaymath}
is the discrete distribution for the post-collision particle
species $i$, given a collision process of type $k$.

Finally, writing
\begin{displaymath}
C_{k} = \sum_i n_{k,i} w_{k,i}(\vv)
\end{displaymath}
defines the post collision velocity distribution $w_{k,i}$, given: a collision has occurred at
pre collision state $(\vr ' , \vv ' , i ')$, then a type of process $k$ has been sampled, then a post collision species
$i$ has been identified. Sampling $(\vv)$ from $w_{k,i}(\vv)$ is carried out in routines VELOCX, VELOEL, VELOEI, VELOPI,
depending on the type $k$ of process identified. In all cases first a transformation into the rest frame of the
interacting background species $i_b'$ is carried out.

\subsection{elastic collisions, VELOEL}

In elastic collisions (EL) colliding particles retain their identity: ($i' = i,  i_b = i_{b'}$). A binary collision
is simulated by first identifying the velocity $\vv_{b'}$ of the pre collision partner from the background medium.
As in all heavy particle collisions the distribution for finding $\vv_{b'}$ is:
\begin{displaymath}
w_{k,i_{b'}}(\vv_{b'}) \sim \sigma_{el}(|\vv'-\vv_{b'}|) |\vv'-\vv_{b'}| f_{i_{b'}}(\vv_{b'})
\end{displaymath}
with $f_{i_{b'}}(\vv_{b'})$ the velocity distribution of background medium $b$. EIRENE typically assumes here
either a drifting Maxwellian in the laboratory frame, or a mono-energetic isotropic distribution in the rest frame of background medium $b$.
Next the scattering angle $\theta$ is found in the center of mass frame, either from the total cross section and the interaction
potential, for from differential cross sections, for other, more approximate assumptions. Finally the post collision velocity
$\vv$ is obtained by transforming back into the laboratory frame.
\subsection{charge exchange, VELOCX}
In this case an exchange of identity is assumed:  $i = i_{b'},   i_b = i'$, which corresponds to a scattering angle
$\theta = \pi$ in the center of mass system. In this special case sampling is most convenient in the
in the rest frame of background medium $b$, rather than in the center of mass frame. This means again first sampling
$\vv_{b'}$ from
\begin{displaymath}
w_{k,i_{b'}}(\vv_{b'}) \sim \sigma_{cx}(|\vv'-\vv_{b'}|) |\vv'-\vv_{b'}| f_{i_{b'}}(\vv_{b'})
\end{displaymath}
and then, after transformation from the rest frame of species $b$ back into the laboratory frame, simply setting $\vv = \vv_{b'}$
\subsection{electron-impact collisions, VELOEI}
\subsection{general heavy-particle-impact collisions, VELOPI }
\subsection{photon processes (emission, absorption, scattering)}
documentation to be written
\newpage
\section{Surface Reflection Models} \label{sec1.3}

Within the terminology introduced in \cref{sec1.1} the effects
of the interaction of neutral particles, photons (and escaping ions) with
surfaces surrounding or inside the computational volume can be
described by additional boundary conditions.
Such kinetic boundary conditions are then often referred to as ``bidirectional reflectance functions''.
It is, however, much
more convenient for Monte Carlo applications to describe this
interaction as just one special type of collision process in the
collision kernel $C$ of equation \cref{eq1.3}, i.e. as
\begin{equation} \label{eq1.14}
C   =   p_w    \cdot   C_w    + \sum_{k}  p_k   C_k   ,
\end{equation}
with $ p_w (\vr )$ = 1, if $ \vr $ is a point
on a surface, and $ p_w (\vr)$ = 0 elsewhere. The
transport kernel $T$ can be modified such that the maximum length
$ l_{max}$ for free flight to the next point of collision is the
distance to the nearest surface along the track of the particle.
The reflection kernel $ C_w$ is further decomposed it into a
kernel $ C_{wf}$ (fast particle reflection), $ C_{wt}$ (thermal
particle re-emission) and $ C_{wa}$ (particle absorption) with
respective probabilities, $ p_f , p_t$  and $ p_ a$, such that
\begin{equation}
C_w   = C_{wf}   +   C_{wt}   +   C_{wa}   =
  p_f   {\cal{C}}_{wf}   +
  p_t   {\cal{C}}_{wt}   +   p_a
{\cal{C}}_{wa}  , \quad p_f + p_t + p_a = 1 . \label{eq1.15}
\end{equation}
Here ${\cal{C}}$ are, as above, the normalized versions of the collision kernels $C$.

We discuss in detail only the fast particle reflection model $
C_{wf}$ since $C_{wt}$ is generally a mono-energetic and cosine
angular distribution or a Maxwellian flux distribution at wall
temperature, and $C_{wa}$ only describes the transition into the
``limbo state" $ x_a$ ``absorbed particle".
\\ It is usual
in neutral gas models to replace $C_{wf}$ with numerical or
analytical fits to moments of $C_{wf}$, e.g., such as the
particle and energy reflection coefficients $ C_P$ and $ C_E$. The
particle - wall interaction is often supplemented by simple
assumptions on the angular distribution (cosine or specular) of
the particles leaving the surface.

One has  for the particle reflection coefficient:
\begin{equation}
C_P   =   \sum_{i}   \int d \vv \   C_{wf} (
\vr ; \vv ' , i ' \rightarrow   \vv
,i) \label{eq1.16}
\end{equation}
Note that $ C_P   =   p_f = c_{wf}$ in the terminology of
\cref{eq1.4}.

The energy reflection coefficient is defined as:
\begin{equation}
C_E    =    \frac{1}{E '}    \cdot \sum_{i}   \int d \vv
\   E    \cdot C_{wf} (\vr ; \vv ' , i '
\rightarrow   \vv ,i)   =
 \frac{1}{E '}     \cdot   {\tilde{C}}_E
\label{eq1.17}
\end{equation}
Here $E = m v^2 / 2$ is the energy of the reflected particle.
Hence $C_E = \bar {E}/E' \cdot C_P$, with $\bar {E}$ denoting the
mean energy of the reflected particles.

\subsection{The Behrisch matrix reflection model} \label{sec1.3.1}
One model, the matrix model due to Behrisch \cite{kn:Behrisch}, is
hardwired in EIRENE (and in many other codes such as AURORA,
BALDUR, TRANSP developed in the early eighties of the last century).
 This model has been used as a standard option in
many benchmarks \cite{kn:Benchmark}. The model is based upon a
table of values for $ C_{P_\perp} (E ')$ and a stochastic matrix $
{\underline{\underline{B}}}_{\perp}$ for the transition $ E '
\rightarrow   E$, for H atoms striking a stainless steel target at
normal incidence. For other incident angles $ \vartheta_{in}$ the
moments are modified in EIRENE by the relations:
\begin{equation}
C_P ( \vartheta_{in} )   =   1   - (   1   -   C_{P_\perp})
 \cdot   \cos^{e1} ( \vartheta_{in} )
\label{eq1.18}
\end{equation}
and
\begin{equation}
C_E ( \vartheta_{in} )   =
  1   -   (   1   -   C_{E_\perp})
 \cdot   \cos^{e2} ( \vartheta_{in} )    .
\label{eq1.19}
\end{equation}
This latter moment equation can be satisfied by modifying the
tabulated stochastic matrix $ {\underline{\underline{B}}}_{\perp}
( E ' \rightarrow   E)$ in many different ways. In EIRENE, we
first sample $ {\cal{E}}$ from $ {\underline{\underline{B}}}
_{\perp}$ (given $ E '$) and then set
\begin{equation}
E   =   E '   -   (   E ' - {\cal{E}}  )
 \cdot   \cos^{e2} ( \vartheta_{in} )    .
\label{eq1.20}
\end{equation}
The linearity of E in $ {\cal{E}}$ gives an expectation value $\bar
{E}$ with $ \bar {E} / E ' \cdot C_P( \vartheta_{in} )  = C_E (
\vartheta_{in} )$, with $C_P $ and $C_E$ as given by equations
\cref{eq1.18} and \cref{eq1.19}, respectively.

To describe the angular distribution of re-emitted particles, conditional
on incident energy and angle,  let
$ \ve_x ,   \ve_y ,   \ve_z$ define
an orthonormal basis at the strike-point of a test-particle on a
surface. Here let $ \ve_x$ be the unit vector parallel
to the outer surface normal and let the incident particle travel
in the local xz-plane. Furthermore, let $ \vartheta_r$ and $
\varphi_r$ be the polar and azimuthal angles sampled from a cosine
distribution around $ - \ve_x$. The speed unit vector $
\vOmega $ of the reflected particle in EIRENE is then
given by
\begin{align}
\Omega_x & =  - \cos( \hat{\vartheta}_r ) f_1 ( \vartheta_{in}
)   +   \sin( \hat{\vartheta}_r )
\sin( \varphi_r)   f_2 ( \vartheta_{in} )     ,   \nonumber\\
\Omega_y & =  \sin( \hat{\vartheta}_r )
\cos( \varphi_r )     ,                           \nonumber\\
\Omega_z & =  \sin( \hat{\vartheta}_r ) \sin( \varphi_r )
f_1( \vartheta_{in} )   + \cos( \hat{\vartheta}_r )   f_2(
\vartheta_{in} )  , \label{eq1.21}
\end{align}
with
\begin{align}
\sin( \hat{\vartheta}_r ) & = f_1( \vartheta
_{in})   \sin( \vartheta_r )     , \nonumber\\
 f_1( \vartheta_{in} )    & = \cos^{e3} ( \vartheta_{in} ) ,
\nonumber\\
 f_2( \vartheta_{in} )    & =
(   1   -   f_1  \cdot f_1 )^{\frac{1}{2}} . \nonumber
\end{align}
In total we have defined three parameters e1, e2, e3. The first
two parameters e1 and e2 account for particle and energy
reflection coefficients that continuously increase with incident
angle (recommended: e1=1, e2=0.5 from comparison with TRIM code
calculations, see below).
\\ The third parameter e3 can be used to take into
account an increasing contribution of specular reflection at near
grazing incidence. e3=0 gives a pure cosine distribution, e3=1
reproduces the distribution given in \cite{kn:Heifetz}. In general
this model describes a fraction of specular versus cosine
reflection, which increases with $ \vartheta_{in}$. And in the
(purely mathematical) limit of $ \vartheta_{in}   =   90^{\circ}$,
it results in pure specular reflection. Furthermore, the parameter $e3$ controls
the speed at which the specular part $a_{spec} (
\vOmega; \vartheta_{in})$ of the angular distribution $a
(\vOmega; \vartheta_{in})   =   a_{spec} ( \vOmega; \vartheta_{in})
  +   a_{cosin} ( \vOmega; \vartheta_{in})$ increases with incident angle
$ \vartheta_{in} $.

A simpler, ``fixed momentum accommodation coefficient model", in which the dependence of this decomposition
on incident angle $\vartheta_{in} $ is absent, can also be activated, e.g.  for testing, (see flag: AINTG
in cards for surface reflection models, \cref{sec2.6}).
\\
The interaction of neutrals and ions with a solid surface for
target-projectile combinations other than H onto stainless steel
are modelled using a reduced energy scaling.

\subsection{TRIM code database reflection models} \label{sec1.3.2}
Complementary to the surface model described above, EIRENE is
coupled to computer generated reflection databases, produced with
the TRIM code (see references \cite{kn:Eckstein} and
\cite{kn:Bateman}). At present a discretized form of the
reflection kernel $C_{wf}$, including all correlations between the
five vector components of $ \vv '$ and $\vv $, is
available for the target-projectile combinations listed in
\cref{tabref}. The first 12 of these combinations
comprised an old default ``TRIM data file" of such combinations in an EIRENE run.
Meanwhile a much extended list,
also far beyond that of \cref{tabref}, is available, see separate
documentation of TRIM code bidirectional reflectance database,
and the target projectile combinations
to be used in any particular run are selected individually
from this larger list (input block 6).

ZNML is a target material flag set for each surface
in the input file, blocks 3a and 3b, and is described below in
block 6b, ``Data for local reflection models". These databases
are, essentially, tables of conditional quantile functions
obtained from large random samples of test particle trajectories
in the solid.
\begin{table}[h]
\caption{\textbf{Surface reflection database (TRIM Code)}}
\label{tabref}
\begin{center}
\begin{tabular} {|r|l|l|r|} \hline
 \textbf{No} & \textbf{Projectile} & \textbf{Target}
 &  \textbf{ZNML}
 \\ \hline \hline

1 & H & Fe & 5626.   \\ \hline

2 & D & Fe & 5626.    \\ \hline

3 & H & C  & 1206.    \\ \hline

4 & D & C  &   1206.    \\ \hline

5 & He &Fe &  5626.    \\ \hline

6 & He & C  &  1206.    \\ \hline

7 & T  & Fe & 5626.    \\ \hline

8 & T & C  &  1206.    \\ \hline

9 & D  & W  & 18474.    \\ \hline

10 & He & W  & 18474.    \\ \hline

11 & H  & W  & 18474.    \\ \hline

12 & T  & W  & 18474.    \\ \hline\hline

13 & D  & Be &  904.    \\ \hline

14 & D  & Mo & 9642.    \\ \hline

15 & Ne & C  & 1206.    \\ \hline

16 & Ne & Be & 904.    \\ \hline

17 & H  & Cu & 6429.    \\ \hline

18 & H  & Mo & 9642.   \\ \hline

19 & T  & Mo &  9642.   \\ \hline

20 & He & Mo &  9642.   \\ \hline
\end{tabular}
\end{center}
\end{table}

We consider these ``database -" reflection models to be the most
complete option available at present and recommend to further
extend the databases to additional projectile - target
combinations, and in particular to sputtering and self-sputtering
data (such as, e.g., C onto C, or Be, etc.).
\\ For an INTOR benchmark case for neutral particle Monte Carlo Codes
\cite{kn:Benchmark} the choice of a particular surface reflection
model had no strong impact on the results. This was shown by
running the EIRENE code with all three models (``Behrisch Matrix",
TRIM database and MARLOWE database (at that time available as a
third option. Meanwhile the MARLOWE database is not available
anymore).

It was found that this had little influence on global parameters
as, e.g., the pumping efficiency for the particular geometry (5.4
\%, 5.3 \% and 5.2 \% respectively.). On the other hand, it can
easily be imagined that in other cases (different plasma
conditions, more detailed geometries) the choice of the reflection
kernel $ C_{wf} $ can be more important for the result. Certainly
in all cases, in which the results are sensitive to the neutral
particle velocity distribution near surfaces, details of the
reflection kernel matter. This is, for example, the case for
computed line shapes of certain electronic transitions ($
H_{\alpha} $ lines etc.), see e.g., reference \cite{kn:Reiter92d}.

\newpage
\section{Recycling surface sources} \label{sec1.4}

It has been found (mid eigthies of last century: within INTOR framework) as a result of benchmark calculations with
several different neutral gas transport Monte Carlo codes (EIRENE, DEGAS, NIMBUS, DDC83) on the
same input data that the main divergence in the results from the
different codes tested was due to different primary source terms
$Q$ in the transport equation \cref{eq1a,eq1b,eq1.6a}, although target fluxes and temperatures had been
prescribed to be identical. Here we describe one particular EIRENE
surface source option in detail. The way the source is modelled by
EIRENE allows  consistency with kinetic boundary
conditions for plasma fluid equations at plasma facing surfaces.
To describe this, we first write the volumetric source $Q$ (particles per unit time and per unit $(\vr,\vv)$-phase space volume)
in terms of a  surface flux $\Gamma_w$ and a surface delta-function. The
surface (wall) flux $\Gamma_w$ thus has dimension: particles per unit time, per velocity space unit, and per unit area:
\begin{equation}
Q_w( \vr , \vv , i )   = \Gamma_w (
\vr, \vv , i )    \cdot   \delta_w ( \pi )     ,
\label{eq1.22}
\end{equation}
where $\pi$ is a coordinate normal to the surface $S_w$:
$\pi(\vr)=0 $, i.e.: $\ve_{\pi} =
\nabla_{\vr}\pi(\vr)$, pointing away from the
computational volume.
We see that $\Gamma_w$ specifies the (ingoing) ``transport" flux $\Phi$ introduced in \cref{trans-flux} at the boundary of
the computational domain.
We will henceforth use the notation $\Gamma$ rather than $Q$, in case of surface sources.
$\Gamma_w$ is a kinetic (microscopic) surface source (flux-)
density distribution, backward directed (into the computational
volume, i.e., it is non-zero only for $v_\pi < 0$).



To convert from a forward directed ion flux distribution
$\Gamma_{Sh,ion}(\vv)$ at the plasma-sheath interface $S_{Sh}$ to a
neutral flux distribution back from the wall surface $S_w$, we
fold $\Gamma_{Sh,ion}$ with a (magnetic plus electrostatic) sheath transmission
kernel $T_{Sh}$ and with the surface reflection (incl. sputtering) kernel $C_w$,
equation \cref{eq1.14}:
 \begin{align}
\Gamma_w ( \vr , \vv  , i ) & = \sum_{i '}
\int\! \int\! \int d^3 r' d^3 v'  d^3
v'' \Gamma_{Sh,ion} ( \vr ' ,
\vv '' , i ' ) \nonumber\\
& \phantom{=} \cdot T_{Sh} (i ' ; \vr ' ,
\vv '' \rightarrow \vr , \vv  ' )
 \cdot  C_w ( \vr ;
\vv ' , i ' \rightarrow \vv , i )
\label{eq1.24}
\end{align}
The summation over $i '$ is performed over all ion species
entering the sheath region.

A typical example for $\Gamma_{Sh,ion}$ is at least sonic parallel plasma flow (Mach number $M_\|\geq 1$ along the
magnetic field lines) entering the magnetic pre-sheath. In this magnetic pre-sheath
the ion trajectories are bend over (accelerated by the electric pre-sheath field) such that at the entrance of the electrostatic sheath $\vr'_{eSh}$
(locate intermediate between magnetic pre-sheath entrance $\vr '$ and target surface $\vr$ at least sound speed is already reached with respect to the surface normal: $M_\pi\geq 1$.

If the magnetic field is directed normal to the target surface, then the magnetic pre-sheath is absent, $\vr' = \vr'_{eSh}$  and, of course
$M_\| = M_\pi \geq 1$.

A commonly adopted simplification is made in plasma fluid descriptions by ignoring the
typically small spatial volume between the plasma sheath interface
and the wall surface. Hence: $T_{Sh}$ factorizes into a 3D delta
function $\delta^3(\vr-\vr')$ and a deterministic transition
kernel for the velocities $\vv'' \rightarrow \vv'$.



Random numbers from $\Gamma_w$ are generated conveniently by
sampling, firstly, $(\vr ' , \vv  '' , i ' ) $
from $\Gamma_{Sh,ion}$, next generating $(\vr ,
\vv ' ) $ from $T_{Sh}$ and then, finally, sampling the weights, velocities and species indices of re-emitted (neutral) particles
$(\vv , i)$ from the conditional distribution $C_w $, which is conditional on one ion (species $i'$) having hit the
target at position $\vr' = \vr$ and with velocity $\vv '$.

A frequently employed assumption for $\Gamma_{Sh,ion}$ is a one sided (only positive velocity components into the sheath)
forward drifting Maxwellian ion flux distribution at the interface between plasma
and electrostatic sheath $S_{Sh}$. In this case the bending of trajectories from parallel (to $\vB$) to
normal (to the surface) sonic flow within the magnetic pre-sheath is already included in $\Gamma_{Sh,ion}$,
e.g., via boundary conditions in a \ac{CFD} model for plasma ion flow, and the action of
operator $T_{Sh}$ is solely the acceleration in the electrostatic sheath.

Alternatively, and in cases when the magnetic field is tilted with respect to the wall surface,
one may define $\Gamma_{Sh,ion}(\vv)$ to be the kinetic transport flux distribution entering the magnetic pre-sheath,
i.e. with a drift velocity of at least ion sound speed in the direction parallel to the magnetic field $\vB$ (rather than normal to the wall).
In this case the action of the magnetic pre-sheath must be accommodated in the sheath kernel $T_{Sh}$, which we might
even consider to split into a magnetic pre-sheath (MP) kernel and an electrostatic sheath (Sh) kernel:
\begin{equation}
T_{Sh} =  T_{mpSh}(\vr_{mpSh}, \vv_{mpSh} \rightarrow \vr_{eSh}, \vv_{eSh}) \cdot  T_{eSh}(\vr_{eSh},\vv_{eSh} \rightarrow \mbox{wall})
\end{equation}
with the subscripts $mpSh$ and $eSh$ labelling the plasma side magnetic pre-sheath entrance and the electrostatic sheath entrance, respectively.

\begin{figure}[h]
\includegraphics[height=9.0cm]{magn_pre_sheath}
\caption{Transition region between fluid plasma (\ac{CFD}) and wall surface, schematic. The plasma flow is specified by \ac{CFD} models
at the entrance surface of the magnetic pre-sheath MP (taken normal to the magnetic field $\vec{B}$). $\Delta \Phi$ is the fraction of the total
sheath potential drop inside the electrostatic sheath.}
\label{magpresheath}
\end{figure}

\subsection{truncated drifting Maxwellian flux at electrostatic sheath entrance}\label{maxw-flux}

One of the standard options in EIRENE, in particular in applications in which EIRENE is coupled iteratively to a plasma fluid code,
is a truncated forward drifting Maxwellian ion flux distribution (options: NEMODS = 6,7, \cref{sec2.7}).
The kinetic plasma ion flux distribution entering the electrostatic sheath is given as
\begin{equation}
\Gamma^0_{Sh,ion}   = \frac{1}{\alpha^2}   \cdot   v_{\pi}   \cdot
  \exp \left ( - \frac{1}{v_{T_i}^2}
(\vv - \vV_d )^2  \right )    , \quad
v_{\pi} > 0 \label{eq1.23}
\end{equation}
with $v_{T_i}  =   \sqrt{2 kT_i/m_i }$, the thermal ion velocity, normalizing constant
$\alpha^2$ and with the drift velocity $\vV_d$. See \cref{sec2.7} for a description of related EIRENE input flags.
The
superscript 0 indicates normalization of the flux to one particle per unit time and unit area:
\begin{equation}
\int_{v_\pi>0} d^{\ 3}v~ \Gamma^0_{Sh,ion}(\vv) = 1
\end{equation}

Consider now the three orthogonal velocity components $v_1,v_2,v_3$ chosen such that $v_3 = v_\pi$, i.e. the $v_3$ velocity component is normal
to the surface (target), and hence also normal to the plasma-sheath interface.
$\Gamma^0_{Sh,ion}(\vv)$ factorizes into the three uni-variate distributions
\begin{equation}
\Gamma^0_{Sh,ion}(\vv)=\Gamma^0_{Sh,ion,1}(v_1)\cdot \Gamma^0_{Sh,ion,2}(v_2) \cdot \Gamma^0_{Sh,ion,\pi}(v_\pi)
\end{equation}
The first two factors here are simply (shifted) Gaussian distributions and sampling random velocities $v_1, v_2$ from them is trivial (e.g.: Box Muller method),
whereas the third factor is a one-sided drifting Maxwellian flux distribution:
\begin{equation}
\Gamma^0_{Sh,ion,\pi}(v_\pi)   = \frac{1}{\alpha^2_\pi}   \cdot   v_{\pi}   \cdot
  \exp \left ( - \frac{1}{v_{T_i}^2}
(v_\pi - V_{d,\pi} )^2  \right )    , \quad
v_{\pi} > 0
\end{equation}
The normalizing factor $\alpha_{\pi}^2$ for this (now uni-variate) distribution of the normal component  $v_\pi$
  reads:
\begin{equation}
\alpha_{\pi}^2   =  \frac12 v_{T_i}^2  \cdot  \exp(-{\cal V}^2_{\pi} )
   g({\cal V}_{\pi} )  ,
\label{eq1.27}
\end{equation}
with ${\cal V}_{\pi} = V_{d,\pi}/   v_{T_i} $ and $
g(x)   =   1   +   \sqrt{\pi} ~  x ~  [1 + \mathrm{erf}(x)] ~  \exp(x^2
)$.

(Note that the dimensionless velocity $\cal V$ normalized by the thermal speed introduced here is related to the Mach number $M$
(normalized by the ion sound speed) by
\begin{equation}\label{normfac}
{\cal V} = \frac{V}{v_{T_i}} = \frac{V}{c_s}~ \frac{\sqrt{Z_i~T_e/T_i+1}}{\sqrt{2}} = M ~\frac{\sqrt{Z_i~T_e/T_i+1}}{\sqrt{2}}
\end{equation}
with $M = V/c_s$ and the (isothermal) ion sound speed $c_s$ defined as: $c_s=\sqrt{(Z_i~T_e+T_i)/m_i}$
$Z_i$ is the ion charge state.
Hence ${\cal V} = M$ only if $Z_i~T_e=T_i$.

It is clear that random sampling the $\pi$-component $v_{\pi}$ of $\vv $ from
$\Gamma^0_{Sh,ion,\pi}$ is not trivial even for this seemingly simplest sheath entrance ion velocity distribution,  because the cumulative
distribution function $F$ of a 1D forward drifting Maxwellian (with
cut-off at $v_{\pi}=0$):
\begin{equation}
F ( v_{\pi} )   = \int_{0}^{v_{\pi}} dv_{\pi} ' \
\Gamma^{0}_{Sh,ion,\pi} ( v_{\pi} ' ) \label{eq1.25}
\end{equation}
cannot be inverted explicitly.


Therefore in the EIRENE code we apply the Monte Carlo  ``biassed source sampling" concept, (module: \texttt{VELOCS.F}).
We sample this component
${\tilde{v} }_{\pi}$ from a ``non-analog" truncated Maxwellian flux distribution
$\tilde\Gamma^0_{Sh,ion,\pi}(\tilde v_{\pi})$ (see \cref{sec1.2.1a} above)
with zero drift in $\pi$-direction but with different ``non-analog" $\pi$- temperature ${\tilde{T}}_{\pi,i},
({\tilde{T}}_{\pi,i}   \ne T_i   , {\tilde{V}} _{d,\pi} = 0 )$.
\begin{equation}
\tilde\Gamma^0_{Sh,ion,\pi}(\tilde v_\pi)   = \frac{1}{\tilde \alpha^2_\pi}   \cdot   \tilde v_{\pi}   \cdot
  \exp \left ( - \frac{1}{v_{\tilde T_{\pi,i}}^2}
(\tilde v_\pi )^2  \right )    , \quad
\tilde v_{\pi} > 0
\end{equation}
Note that although this non-analog
distribution clearly violates any ``Bohm criterion", it is, taken together with the weight correction factors described below, unbiased and
a fully legitimate sampling distribution because it fulfils the ``Radon-Nikodym condition''
\cref{eq1.14a}.
Sampling from this non-analog distribution can
be done by setting, e.g.,
\begin{equation}
\tilde{v}_{\pi}   = {\tilde{\alpha}}_{\pi}   \cdot   \sqrt{ -\ln(
\xi )  \cdot 2 } = v_{\tilde T_{\pi,i}} \cdot   \sqrt{ -\ln(
\xi ) }  \quad  \mbox{since: \ }  \tilde{\alpha}_{\pi}   =\frac{1}{\sqrt{2}}
v_{\tilde T_{\pi,i}} \quad \mbox{for \ } \tilde{V}_{d,\pi} = 0
\label{eq1.26}
\end{equation}
$\xi$  being a uniformly distributed random number between 0 and
1.

\smallskip

In order to find the weight correction factor associated with this non-analog
procedure, we first note that from \cref{normfac} we have
\begin{displaymath}
{\tilde{\cal V}}_{\pi} = 0 \rightarrow g({\tilde{\cal V}}_{\pi}) =1, \mbox{and}  \quad \tilde{\alpha}_{\pi}   = \frac{1}{\sqrt{2}}v_{\tilde{T_i}}
\end{displaymath}
This gives the following result for the (multiplicative) weight
correction factor $w ( \tilde{v}_{\pi} )$, see equation \cref{eq1.13a}, to be applied for each sampled non-analog velocity component
$\tilde v_\pi$:
\begin{equation}
w ( \tilde{v}_{\pi} ) = \frac{\Gamma^0_{Sh,ion,\pi} (
{\tilde{v}}_{\pi} )}{\tilde{\Gamma}^0_{Sh,ion,\pi} ( \tilde{v}_{\pi}
)}
 =  \frac{\tilde{T}_{\pi,i}}{T_i}
\exp \left (  \frac{2 \tilde{v}_{\pi}  V_{d, \pi} }{v_{Ti}^2}    -
\frac { {\tilde{v}_{\pi}}^2}{v_{T_i}^2} \left ( 1   -
\frac{T_i}{ {\tilde{T}_{\pi,i}} }  \right ) \right )
 \cdot  \frac {1}{g({\cal V}_{\pi})}
\label{eq1.28}
\end{equation}
This non-analog sampling of the source term introduces additional
statistical fluctuations in the results, since the initial test
particle weights $w$ fluctuate with variance
\begin{equation}
\sigma^2(w) = \int_{0}^{\infty} \! d \tilde{v}_{\pi} \ w^2 (
{\tilde{v}}_{\pi} )  \cdot \tilde{\Gamma}^0_{Sh,ion} (
\tilde{v}_{\pi} )
 - 1 = g \left ( \frac{2{\cal V}_{\pi}}{\sqrt{b}} \right ) /
g({\cal V}_{\pi})^2
 \cdot  \frac{1}{(2-b)   b}   -   1
\label{eq1.29}
\end{equation}
where $b   =   2   -  T / {\tilde{T}}$. We use the
ansatz $b = C  \cdot   {\cal V}_{\pi} + 1$ and compute numerically that
value of C which minimizes $\sigma^2$ in the range $0 \le {\cal V}_{\pi}
\le 1$. We find C = 0.6026. This means that the ratio of the
``real" temperature $T_i$ and the non-analog temperature $\tilde{T}_{\pi,i}$
should be chosen to be $ T_{i} / {\tilde{T}_{\pi,i}}   =   1 - 0.6026  \cdot
{\cal V}_{\pi}$. This choice optimizes the statistical performance in the given Mach number range. In
principle, of course, any other $\tilde{T_{\pi}} $ could have been used
without introducing bias into
the algorithm.\\
The mean energy of ions generated by this sampling procedure is
(note: the distribution $\Gamma^0$ is  already normalized to flux
one):
\begin{align}\nonumber
\overline{E}({\cal V}_{\pi},{\cal V}_{\sigma})   &=
 \frac{m_i}{2}    \cdot   \int^\infty_{-\infty} \!
\int^\infty_{-\infty} \! \int^\infty_{0} \!
 d^3v' \
 {v'}^2    \cdot
\Gamma^0_{Sh,ion} (\vv ' ) = \\
&= \frac{m_i}{2}    \cdot   \int^\infty_{-\infty} \!dv_1
\int^\infty_{-\infty} \!dv_2 \int^\infty_{0} \!
   dv_\pi \
 (v_1^2+v_2^2+v_\pi^2)    \cdot
\Gamma^0_{Sh,ion} (\vv ' ) \label{eq1.30} \\\nonumber
&= T_i    \cdot   \gamma_E    ,
\end{align}
with the half sided integration for the $\pi$-component of
$\vv$, and:
 \begin{equation}
\gamma_E   =    2 + {\cal V}_\sigma^2    + \left[{\cal V}_\pi^2    +   0.5
\cdot \frac{g({\cal V}_{\pi} )  - 1}{g({\cal V}_{\pi} )}\right]       , \label{eq1.31}
\end{equation}
Here we have decomposed the drift velocity $\vV_d$ into a
normal (to the surface) $\pi$-component $\vV_{d,\pi}$ and a $\sigma$-component $\vV_{d,\sigma}$ parallel to
the wall surface.
\begin{equation}
{\cal V}_{\sigma}   = \frac{ | \vV_d   -   \vV_{d,
\pi} | } {  v_{T_i}} = \frac{ | \vV_{d,\sigma} | } {    v_{T_i}}. \label{eq1.32}
\end{equation}
Note that for ${\cal V}_{\pi}$ = 0 \cref{eq1.31} reduces to the expected formula
$\gamma_E   =   2   + {{\cal V}_{\sigma}}^2$  for a half sided Maxwellian flux with only drift components parallel to the surface.

Furthermore, if ${\cal V}_{\sigma}
= 0$ (normal incidence of the mean flow) and ${\cal V}_{\pi} \rightarrow 1$ (i.e. for the Bohm sheath
condition $M_{\pi} \rightarrow 1$ in case of $Z_i~T_e = T_i$, such that $cs = v_{T_i}$), then $\gamma_E$ for the forward sided
fluxes approaches the value for total net heat fluxes from (full) drifting Maxwellian distributions:
$\gamma_E = 2.449 + {{\cal V}_{\pi}}^2 \simeq \frac {5}{2} +
{{\cal V}_{\pi}}^2$, as it must.

We conclude that the above scheme is both an accurate and
efficient random number generator for truncated shifted Maxwellian
flux distributions of particles incident onto or emitted from
surfaces. Of course, after having generated $
\tilde{\vv}  $ with the weight $ w ({\tilde{v} }_{\pi}
)$ and before reflecting the particle via the kernel $C_w $, the
acceleration of an ion in a sheath potential  in direction $\pi$
(normal to the target surface), as formalized in the sheath
transmission kernel $T_{Sh}$ mentioned above in Eq. \cref{eq1.24}, can easily be
included.

\subsection{The electrostatic sheath} \label{sec1.4.1}


The electrostatic sheath potential in front of a target surface is
derived from assuming a given net electrical current $j_{pl}$ and
a known ion particle current (ion species $i$) $\Gamma^+_i$ (i.e., $j^+_i = q_i
\Gamma^+_i$ electrical ion current entering a plasma-target boundary layer (the ``sheath"), with $q_i=eZ_i$ being the ion charge, $e$ the elementary charge and $Z_i$ the ion charge number)
flowing over the target sheath edge. Both these currents $j_{pl}, \Gamma^+_i$, and the
electron current

$j^-=q_e\Gamma^- = -e\Gamma^- = j_{pl} - \sum_i j^+_i = j_{pl} -  j^+$,
(with $j^+$ denoting the total, sum over all ion species, ion current)
are taken to
flow parallel to the $\vB$-field at the plasma-sheath edge (or, more generally: to form (different) angles $\Psi_{\Gamma^+,j_{pl}} > 0$ degrees, respectively,  with the target surface normal $\pi$).

Often, but not necessarily, it is assumed: $j_{pl} = j^+ + j^- =
0$, i.e., locally ambi-polar flow conditions.

We start with a Maxwell-Boltzmann density distribution $n_e(x)$
for the electrons, with electron temperature $T_e$, in the
electron retarding electric field with potential $\Phi(x)$. Here $x$
is a coordinate parallel to the current flow, e.g. parallel to the magnetic field $\vB$ in case of absence of drift flows. The
electron current density along the field is then (after reduction,
possibly, due to secondary electron emission):
\begin{displaymath}
j^-_T = -e \cdot 1/4 \cdot n_{e,T} \cdot v_e^{av} \cdot
(1-\gamma_{s.e.e.}) = -e \cdot 1/4  \cdot n_{e,S} \cdot \exp[e
\Delta\Phi_{S,T}/kT_e] \cdot \sqrt{\frac{8 kT_e}{\pi m_e}} \cdot
(1-\gamma_{s.e.e.})
\end{displaymath}
Here $\gamma_{s.e.e.}$ is the secondary electron emission
coefficient and $v_e^{av}$ is the average electron velocity. The subscripts $S$ and $T$ stand for the plasma-sheath
interface and for the target surface, respectively and $
\Delta\Phi_{S,T}$ is the (typically negative) voltage drop from
the plasma sheath edge to the target.

In order to determine the sheath voltage drop $\Delta\Phi_{S,T}$
we equate the component normal to the target of this current
density with the corresponding component of the net ion current
density:
\begin{displaymath}
\left(-j^-_T\right)_{\pi} =-j^-_T \cdot \cos(\Psi) =
(j^+_T - j_{pl}) \cdot \cos(\Psi) = \left(j^+_T - j_{pl}\right)_{\pi}
\end{displaymath}
Hence the cosine of the angle $\Psi$ of direction of current flow
(e.g. of field line inclination
against the target surface normal, if the current flows parallel to $\vB$) cancels out, as does the elementary
charge $e$. The ion electrical current to the target $j^+_T$ ($ = (j^+_S)_{\pi}$, no ion
sources within the sheath region) is given as
\begin{displaymath}
j^+_T = e \ \sum_i Z_i \cdot n_{i,S} \cdot {{V}_{\pi,i,S}}
 = e \ \sum_i Z_i \cdot n_{i,S} \cdot M_{\pi,i,S} \cdot c_{i,S}.
\end{displaymath}
The sum is over all ion species, $Z_i$ is the charge state of
species $i$, $n_{i,S}$ is the corresponding ion density and
${{V}_{\pi,i,S}}$ the ion flow
velocity component normal to the target, of ion species $i$, subscript $S$ indicated that these values are taken at the sheath entrance.
$M_{\pi,i}$ is the Mach number of the corresponding current component, and $c_i$
is the ion acoustic speed for species $i$, again taken at sheath entrance.

Hence, combining the three previous relations, the sheath potential difference $\Delta\Phi$
in terms of
given values (at sheath entrance) for $ j_{pl},  T_e, \gamma_{s.e.e.} $ and ions flows $\Gamma_i$,
reads (omitting the subscript $\pi$ for the components normal to the target in these quantities):

\begin{align} \label{sheath1}
 e \Delta\Phi/kT_e & =
\ln\left[\frac{\sqrt{2\pi}}{(1-\gamma_{s.e.e.})}\left(\sum_i
\left( Z_i \frac{n_{i,S}}{n_{e,S}}\cdot {V_{i,S}}
\sqrt{\frac{m_e}{kT_e}}\right)-\frac{j_{pl}}{e \cdot n_{e,S} \cdot
\sqrt{\frac{kT_e}{m_e}}}\right)\right] \\
& = \label{sheath2}
\ln\left\lbrack\frac{\sqrt{2\pi}}{(1-\gamma_{s.e.e.})}\left(\sum_i
\left( Z_i \frac{n_{i,S}}{n_{e,S}} \cdot M_{i,S} \cdot c_{i,S}
\sqrt{\frac{m_e}{kT_e}}\right)-\frac{j_{pl}}{e \cdot n_{e,S} \cdot
\sqrt{\frac{kT_e}{m_e}}}\right)\right\rbrack
\end{align}
% \end{subequations}
This expression \cref{sheath1} is evaluated by the function SHEATH
of the EIRENE code for the sheath potential drop $\Delta\Phi$, taking local plasma data for $T_e$ and $n_i$
and ${V_{\pi,i}}$  (index $i$ runs over all ion species comprising the ionic target surface flux. The local electron density $n_e$ at sheath entrance  follows then  from quasi-neutrality.

All these values may have a spatial dependency and are taken at a (random sampled, according to target outflow plasma flux distributions)
point at a boundary surface of a \ac{CFD} plasma code, e.g., at the divertor target boundary
plate, which, strictly, is at the sheath entrance $S$ in front of the target. The energy gain due to sheath acceleration for a particular ion $i_0$
is then $Z_{i_0} e \Delta\Phi$.

Special forms of \cref{sheath2} are often quoted in the
literature, e.g. for single ion fluid cases. Writing for the ion acoustic speed $c_i$ explicitly
\begin{displaymath}
c_i = \sqrt{\left(\frac{\gamma kT_i + Z_i kT_e}{m_i}\right)}
\end{displaymath}
($\gamma$ denoting the adiabatic coefficient) and hence:
\begin{equation}\label{insert}
c_i \cdot \sqrt{ \frac{m_e}{kT_e} } = \left[ \left( \frac{\gamma
kT_i}{kT_e} + Z_i \right) \cdot \frac{m_e}{m_i} \right]^{1/2}
\end{equation}
we recover well known expressions for simpler cases. Consider,
e.g., a single ion fluid case, $n_e = Z_i n_i = n$,  and
$M_i = 1$ (Bohm sheath condition for isothermal flow, $\gamma =
1$), \cref{sheath2} reduces to:
\begin{align} \label{sheath3}
 e \Delta\Phi/kT_e & =
\ln\left[\frac{\sqrt{2\pi}}{(1-\gamma_{s.e.e.})}
\left({c_{i,S}}
\sqrt{\frac{m_e}{kT_e}}\right)-\frac{j_{pl}}{e \cdot n \cdot
\sqrt{\frac{kT_e}{m_e}}}\right] \\
& = 
\ln\left[\frac{\sqrt{2\pi}}{(1-\gamma_{s.e.e.})}
\left[ \left( \frac{
kT_i}{kT_e} + Z_i \right) \cdot \frac{m_e}{m_i} \right]^{1/2}\left(1-\frac{j_{pl}}{e \cdot n \cdot
c_{i,S}}\right)\right]
\label{sheath4}
\end{align}
the latter equality \cref{sheath4} follows after insertion of \cref{insert} into the former.



If we let, furthermore, the net electrical current be zero: $j_{pl} = 0$ and take
zero secondary electron emission $\gamma_{s.e.e} = 0$, we then find the classical result:
\begin{displaymath}
 e \Delta\Phi/kT_e = \frac{1}{2} \ln{ \left[ \frac {2\pi
m_e}{m_i} \left( \frac{T_i}{T_e} + Z_i \right) \right]}
\end{displaymath}
Hence, e.g.,  for a pure hydrogen plasma ($m_e/m_i \approx 1/1836$), and for $T_e \approx
T_i$, we have $e \Delta\Phi/kT_e \approx -2.5 $ and $ \approx
-2.84$ for a pure deuteron plasma. For a pure helium plasma (fully ionized,
$Z_i=2$) this sheath factor becomes $\approx -2.98$.


\subsection{The magnetic pre-sheath} \label{sec1.4.2}
In most situations when EIRENE is applied to plasma surface interaction studies the ion-plasma flow velocity $\vV_i$
is predominantly parallel to a $\vB$ field: $V_{\parallel,i} \approx |\vV_i|$, and it is this parallel ion flow velocity,
which is usually available from \ac{CFD} plasma simulations, via boundary conditions specified for the (parallel component of the) momentum balance equation
at the entrance of the magnetic pre-sheath ($mpSh$) rather than at the entrance of the electrostatic sheath $eSh$.

The sheath model of EIRENE then is to be supplemented by a model \textbf{PRE-SHEATH} for the transition of known currents and $\vV_{i,mpSh}$  to $\vV_{i,eSh}$, i.e.
in particular to derive the currents $j_{\pi,eSh}$ and flow components $V_{\pi,i,eSh}$ entering the electrostatic sheath from $\vj_{mpSh}$ and $\vV_{i,mpSh}$

to be written....

\newpage
\section{Combinatorial description of geometry} \label{sec1.5}
EIRENE is a Monte Carlo Code. Defining geometry in a Monte Carlo code is complex and laborious.


Of the three typical options for definition of the target geometry in Monte Carlo Codes:
\begin{itemize}
\item {}
    ``\ac{CSG}",
\item
    ``\ac{BREP}"
\item
     ``Voxel Geometry"
\end{itemize}
EIRENE uses a combination of the latter two concepts.

\medskip
In \ac{CSG} a solid object is ``constructed" from the intersection, union or other boolean operations of several ``half-spaces".
Each half space in specified by a relation $F(x,y,z)\le 0$, typically $F$ is a low order polynomial.
For example the neutronics \ac{MCNP} code used e.g. for ITER nuclear safety studies, is based on \ac{CSG}.

\medskip
In \ac{BREP} geometry, surfaces and volumes are build up from constituent parts.
These parts may be related with algebraic functions (elementary \acp{BREP})
or complex functions such as Bezier functions or B-splines (advanced \acp{BREP}).
Most \ac{CAD} software uses \acp{BREP} in their native definition of geometry.

In the EIRENE \ac{BREP} options surfaces are build from (bounded) algebraic functions.
These surfaces are referred to as ``additional surfaces", see \cref{sec2.3b}.
EIRENE uses only algebraic functions, up to second order
(some  extensions to raise the order to fourth order have been started recently but are not ready to use).

\medskip
In ``Voxel geometry", a 2D or a 3D object is constructed from identically shaped volume elements (``voxels"). In EIRENE these may be triangles
 or quadrangles in 2D, or tetrahedra in 3D. New options to accommodate trilinear hexaedra as 3D voxels are currently being developed in connection with 3D edge codes such as EMC3. The bounding surfaces of these voxels are referred to as
``standard grid surfaces" in EIRENE, see \cref{sec2.2}. These volume discretisation meshes (VOXELS) can be read into EIRENE
from external files. For example the mesh in B2 \ac{CFD} plasma fluid code is read into EIRENE directly from the \ac{CFD} grid generator in
coupled B2-EIRENE applications.
Wall meshes are build from these by assigning
a surface reflection model to some of the (voxel-) surfaces (\cref{sec2.3a}), whereas internal (grid-) surfaces remain transparent for flow of test-particles.

\textbf{Note added in 2008:}

An attempt has been made in 2008 to use \ac{CAD} files directly for 3D wall surfaces.
This was done via the already existing \acs{CAD}-\acs{MCNP} interface. We had then tried to turn the \ac{CSG} input
for \ac{MCNP} (see above) into a \ac{BREP} form for EIRENE.  In principle this works. But in practice the
de-tour via the \ac{CSG} concept leads to prohibitively complex \ac{BREP} input for EIRENE.

A direct \acs{CAD}-\acs{BREP} interface seems now much more efficient to us and more natural anyway,
and has meanwhile (2011) been established by an intermediate code step (the ANSYS code), which first converts
\ac{CAD} output into a triangulation of surface elements, and then each triangle (in 3D space) is taken as ``additional surface'' (\ac{BREP} concept)
into EIRENE. Very large sets of triangles required optimisation procedures, which have  also been
implemented (in EIRENE module timea.f) on the basis of OCTREE concepts adopted from ray-tracing procedures used in image processing.


\medskip
A complete set of surfaces in an EIRENE geometry may consist of a combination of both: standard surfaces (hence: defining standard grid cells, (voxels), and additional surfaces, defining a general \ac{BREP} discretisation (additional cells).  A more complicated situation arises when additional surfaces intersect
standard grid voxels.
\begin{figure}[h]
\begin{minipage}{0.48\textwidth}
\includegraphics[width=\textwidth, angle=+90]{ALT-II}
\caption{\label{label1l}Early example a combination of \ac{BREP} and VOXEL discretisation, around 1985
for ALT-II pumplimiter studies at \ac{TEXTOR}. Discretisation inside the \ac{TEXTOR} tokamak vessel (torus): VOXEL, supplemented
by about 100 (\acs{BREP}) first and second order surfaces for ALT-II scoops, pumping ducts and inner structures such as probe supports,
getter pumps, etc..
}
\end{minipage}\hfill%
\begin{minipage}{0.48\textwidth}
\includegraphics[width=\textwidth,angle=-0]{alt_imp}
\caption{%
\label{label1r}same as left Figure, but \ac{BREP} part only, and EIRENE re-run with more modern 3D visualisation graphics program, 2009
\vspace{1.2cm}}
\end{minipage}
\end{figure}




If these surfaces are transparent, they can be used to conveniently ``measure" fluxes at any position.
But if they are reflecting, or absorbing, e.g. to define a complex shaped bounding surface in an otherwise simple grid, then
some care is needed because the automatic cell volume evaluation, which is available in EIRENE for standard grid cells (voxels),
may not correspond anymore to the cell volume accessible for test particles. Cell averaged quantities (divided by cell volume) may then not necessarily
be scaled properly. For a possible work-around see NLERG option in input block 1, \cref{sec2.1}.

From the above it follows than in an EIRENE model
surface elements building up the wall can be defined as bounded algebraic surfaces, with triangles as special case,
i.e. as ``additional surfaces, \ac{BREP}", or can be bounding surfaces of the voxels, if the ``grid of voxels" extends up to the real wall, or a combination of both.

\medskip
As can be seen from the flowcharts, all the geometrical
calculations needed for tracing test flights and for estimating
volume- and surface averaged tallies are compiled in a ``geometry
block", which can easily be exchanged.
In most cases (e.g., all
installations of EIRENE outside FZ-J\"ulich) we load the most
complete, fully 3D block ``GEO3D", which allows fully 3 dimensional spatial
resolution. Lower dimensional (0D, 1D or 2D) are contained as special cases, e.g.
by choosing only grid surfaces with proper symmetry.


There is, in general, not a very large CPU benefit from
running EIRENE with simpler specialized and geometry optimized blocks such as GEO2D or GEO1D for 2D
or 1D models, respectively, because the saving in CPU costs in most
cases is only around 10 - 20 \%, so these options have not been supported since the late eighties of the last century anymore.

\paragraph{The geometry module of EIRENE}\hfill

The following operations have to be executed by the block:
\\Suppose the three dimensional
computational volume is discretized by a mesh of $n$ surfaces
\mbox{$S_i, i=1,2,...,n$ .} Each of these surfaces is defined by a
set of $\Lambda_i$ equations (unbounded surfaces):
\begin{equation}
f_{\lambda}^i (x,y,z)   =   0 { },    \quad \quad \lambda   =
1,2,..., \Lambda_i { }, \label{eq1.34}
\end{equation}
and, optional, for each unbounded surface $\lambda$  a subset of $ \Pi_{\lambda} $
inequalities (e.g. for turning the unbounded surfaces into finite segments):
\begin{equation}
g_{\lambda , \pi }^i (x,y,z)  \le   0 { },    \quad \quad \pi   =
1,2,..., \Pi_{\lambda} { }. \label{eq1.35}
\end{equation}

This means that for all unbounded surfaces $S_{i,\lambda}$ only the points, which fulfill all corresponding
inequalities \cref{eq1.35} are regarded by EIRENE as valid parts of the finite surface element.

Assume a test particle at position $ \vr_0 $ =
(X0,Y0,Z0) moving with speed unit vector $ \vOmega_0 $
= (VELX,VELY,VELZ). The block then has to provide the nearest
intersection $ \vr_1$ of the straight line (``sample trajectory")
\begin{equation}\nonumber
G(t):
 \vr_0   +   t   \vOmega_0
\end{equation}
amongst all possible intersections with the surfaces $S_i$. Moving
the ``particle" from $\vr_0$ to $\vr_1$ and
then repeating the block until the total track length $l$ as
sampled from the transport kernel $T(l)$ \cref{eq1.5} is reached,
allows to generate histories with all information available along
the track to evaluate the estimators mentioned in
\cref{sec1.1}.
\\
In EIRENE $f$ and $g$ are general 2nd order (with linear surfaces
as special case) algebraic equations of the form
\begin{align}
h(x,y,z) = a_0   +   a_1  \cdot  x   + a_2   \cdot  y   +
a_3  \cdot  z  & + a_4  \cdot  x^2   +   a_5  \cdot  y^2   +
a_6  \cdot  z^2   +  \nonumber\\
               & + a_7  \cdot  xy   +   a_8  \cdot  xz   +
a_9  \cdot  yz \label{eq1.36}
\end{align}
Hence in the geometrical block nothing else but second order
equations
\begin{equation}\nonumber
t^2   +   p \cdot  t   +   q   =   0
\end{equation}
have to be solved
accurately and as efficiently as possible. We note in passing that
each of this surfaces may be transparent, purely absorbing, semi
transparent or (partially) reflecting. Surface sources may be
defined on each surface. These surfaces may be ``invisible" from
one side and transparent, absorbing or reflecting from the other
and each surface can ``operate switches" if intersected by a test
flight, such as abandoning the evaluation of atomic or molecular
reaction rates (entrance into vacuum regions) along the flight,
abandoning evaluation of volume averaged responses (e.g., in some
uninteresting regions) or abandoning geometrical calculations (for
those surfaces, for which it can be known a priori that they will
not intersect a flight path originating at a particular surface or
in a particular cell).

\newpage
\section{time dependent mode of operation} \label{sec1.6}
The section is currently rewritten. Please refer to the
J.Nucl.Mater. paper: D.Reiter et al., Vol 220-222 (1995), 987-992.

Details of options available in EIRENE for time dependent problems are also given in \cref{sec2.13}.

The following issues have to be considered in case of time-dependent coupling schemes to external plasma fluid solvers,
see \cref{flow03}:

The coupling is explicit, hence \ac{CFL} type conditions for the
time-step have to be obeyed. Here this means: the time-step must
be shorter than the time between re-ionization and subsequent
neutralization (in the volume or at the targets), because the fate
of a particle during its intermediate ionized phase is not
included in the explicit Monte Carlo procedure.

The noise level for collision estimators increases with smaller
time-steps. Hence: avoid all collision estimators in t-dependent
mode.

Please contact us for an update on this situation, if you wish to
use that option.
\newpage
\section{non-linear effects: coupling to plasma fluid models} \label{sec1.7}
The section is currently rewritten. Please refer to the report
J\"ul-2872, Feb 1994, by G.Maddi\-son and D.Reiter for details (ref.
\cite{kn:Maddison94}).

The back-reaction of the host medium (plasma) on the neutral (test) particle kinetics renders
the combined plasma-neutral transport problem non-linear already for the Monte Carlo solution of the ``linear" Boltzmann equation.
This non-linearity is not in the collision integral, but in the dependence of the collision parameters
(collision rates, mean free pathes, ...) on the plasma parameters. It can be dealt with by iterating a solver for the host medium (plasma) and for the kinetic equation (neutrals).
Such a \ac{CFD} solver for 2D plasma flows is e.g. the B2-``BRAAMS" code, and its iterative coupling to the Monte Carlo kinetic code
EIRENE was developed in the late eighties and early nineties of the last century, under NET-KFA contracts, involving three EFDA associates:
the EIRENE group at FZ J\"ulich (formerly KFA), the CCFE (formerly AEA Culham) edge modelling group as well as ERM Brussels.


The historically oldest option is the ``semi implicit" scheme, see \cref{flow02}, because of
the severe CPU time constrains in the late eighties of the last century, when B2-EIRENE was first
developed, see Reiter et al.,  (\cite{kn:Net}).

Later mostly the simpler ``explicit scheme" , \cref{flow01} was employed in applications, which, although probably being much slower
in ``code turn-around time", did require less attention by the person running the code, i.e. seemed to be
numerically more robust, in particular in cold, recombining divertor plasma conditions.

Since about 2009 also available (but apparently not used in 3rd party applications) is a combination of both methods:
The selection between explicit and semi-implicit iteration is made possible individually for each ``stratum",
such that now some parts of the problem can be run semi-implicitly, while others (notably the volume recombination
neutral gas source) can remain explicit.

All these coupling schemes treat the kinetic part of the model in quasi-steady-state approximation with the plasma state.
This is justified as long as the gas dynamics is fast compared to the plasma dynamics, the latter being roughly specified by the
time step $\Delta t$ used in the \ac{CFD} plasma simulation.

In particular in high density, low temperature detached divertor plasmas this ordering may not necessarily be correct any longer, and a
time dependent mode of operation of EIRENE for the coupling may be required due to \ac{CFL} constrains.

The full exploration of this option,
which was implemented into EIRENE in the mid nineties of the last century for \ac{ELM} simulations, has not yet happened  with respect to
overall coupled code stability and ``turn around code run time". The \cref{flow03} illustrates this mode of operation.








\newpage

\begin{figure}[h]
\includegraphics[width=\textwidth]{flowchart9}
\caption{Explicit coupling scheme between EIRENE and a \ac{CFD} code (for plasma flow):
After each plasma code cycle (or time step) a full new Monte Carlo run is carried out
\label{flow01}}
\end{figure}

%{\center
%{\Huge "explicit coupling"}

%\vspace{2truecm}

%\unitlength1 true cm
%\begin{picture}(16,9) \thicklines
%
%\put(7.0,6.5){\framebox(2,1.5)}
%\put(7.2,7.5){\makebox(0,0)[tl]{LINDA}}
%
%\put(7.0,7.25) {\line(-1,0){3.5}}
%\put(3.5,7.25) {\vector(0,-1){1.5}}
%
%\put(9.0,7.25) {\line(1,0){4.0}}
%\put(13.0,7.25) {\vector(0,-1){4.5}}
%
%\put(2.5,4.25){\framebox(2,1.5)}
%\put(2.7,5.25){\makebox(0,0)[tl]{INTER0}}
%
%\put(3.5,4.25) {\vector(0,-1){1.5}}
%
%\put(2.0,0.75){\framebox(3,2.0)}
%\put(2.2,2.25){\makebox(0,0)[tl]{EIRENE}}
%\put(2.2,1.75){\makebox(0,0)[tl]{nE Histories}}
%
%\put(11.0,0.75){\framebox(4,2.0)}
%\put(11.7,2.25){\makebox(0,0)[tl]{BRAAMS}}
%\put(11.7,1.75){\makebox(0,0)[tl]{nB timesteps}}
%
%\put(7.0,2.0){\framebox(2,1.5)}
%\put(7.2,3.0){\makebox(0,0)[tl]{INTER1}}
%
%\put(7.0,0.0){\framebox(2,1.5)}
%\put(7.2,1.0){\makebox(0,0)[tl]{INTER2}}
%
%\put(7.0,2.25) {\vector(-1,0){2.0}}
%\put(5.0,2.25) {\line(-1,0){0.1}}
%\put(4.9,2.25) {\line(0,-1){1.0}}
%\put(4.9,1.25) {\vector(1,0){2.1}}
%
%\put(11.1,1.25) {\line(0,1){1.0}}
%\put(11.1,2.25) {\vector(-1,0){2.1}}
%\put(9.0,1.25) {\vector(1,0){2.0}}
%\put(11.0,1.25) {\line(1,0){0.1}}
%\end{picture}
%}
\newpage
%{\center

\begin{figure}[h]
\includegraphics[width=\textwidth]{flowchart7}
\caption{``semi implicit" coupling scheme between EIRENE and a \ac{CFD} code (e.g. BRAAMS B2 for plasma flow), involving internal
re-iteration on sources terms without new Monte Carlo trajectories, much less frequent Monte Carlo runs than in ``explicit" mode of operation,
but unresolved issues still re volume recombination contribution.  LINDA is a grid generating tool providing a numerical discretisation mesh which is common for B2 and EIRENE  \label{flow02}}
\end{figure}
%{\Huge "implicit coupling"}


%\vspace{2truecm}

%\unitlength1 true cm
%\begin{picture}(16,12) \thicklines
%
%\put(7.0,10.5){\framebox(2,1.5)}
%\put(7.2,11.5){\makebox(0,0)[tl]{LINDA}}
%
%\put(7.0,11.25) {\line(-1,0){5.75}}
%\put(1.25,11.25) {\vector(0,-1){1.5}}
%
%\put(9.0,11.25) {\line(1,0){4.5}}
%\put(13.5,11.25) {\vector(0,-1){5.75}}
%
%\put(0.25,8.25){\framebox(2,1.5)}
%\put(0.45,9.25){\makebox(0,0)[tl]{INTER0}}
%
%\put(1.25,8.25) {\vector(0,-1){2.75}}
%
%\put(0.0,2.0){\framebox(5,3.5)}
%\put(0.2,5.0){\makebox(0,0)[tl]{EIRENE}}
%\put(2.5,3.5) {\line(0,1){2.0}}
%\put(2.7,5.0){\makebox(0,0)[tl]{EIRENE}}
%\put(2.7,4.5){\makebox(0,0)[tl]{A+M data}}
%\put(0.0,3.5) {\line(1,0){5.0}}
%\put(1.3,3.0){\makebox(0,0)[tl]{nE Histories}}
%
%\put(11.0,2.0){\framebox(5,3.5)}
%\put(11.7,5.0){\makebox(0,0)[tl]{BRAAMS}}
%\put(11.7,4.5){\makebox(0,0)[tl]{(each time-step)}}
%\put(11.0,3.5) {\line(1,0){5.0}}
%\put(11.7,3.0){\makebox(0,0)[tl]{nB time-steps}}
%
%\put(7.0,6.0){\framebox(2,1.5)}
%\put(7.2,7.0){\makebox(0,0)[tl]{INTER3}}
%
%\put(7.0,4.0){\framebox(2,1.5)}
%\put(7.2,5.0){\makebox(0,0)[tl]{INTER4}}
%
%\put(7.1,4.75) {\line(0,1){2.0}}
%\put(7.0,4.75) {\line(1,0){0.1}}
%\put(7.1,6.75) {\line(-1,0){2.2}}
%\put(4.9,6.75) {\vector(0,-1){1.25}}
%\put(4.9,6.75) {\line(0,-1){2.0}}
%\put(4.9,4.75) {\vector(1,0){2.1}}
%
%\put(11.1,4.75) {\line(0,1){2.0}}
%\put(11.1,6.75) {\vector(-1,0){2.1}}
%\put(9.0,6.75) {\line(-1,0){0.1}}
%\put(8.9,6.75) {\line(0,-1){2.0}}
%\put(8.9,4.75) {\vector(1,0){2.1}}
%\put(11.0,4.75) {\line(1,0){0.1}}
%
%\put(7.0,2.0){\framebox(2,1.5)}
%\put(7.2,3.0){\makebox(0,0)[tl]{INTER1}}
%
%\put(7.0,0.0){\framebox(2,1.5)}
%\put(7.2,1.0){\makebox(0,0)[tl]{INTER2}}
%
%\put(7.1,0.75) {\line(0,1){2.0}}
%\put(7.1,2.75) {\vector(-1,0){2.1}}
%\put(4.9,2.75) {\line(1,0){0.1}}
%\put(4.9,2.75) {\line(0,-1){2.0}}
%\put(4.9,0.75) {\vector(1,0){2.1}}
%\put(7.0,0.75) {\line(1,0){0.1}}
%
%\put(8.9,2.75) {\line(0,-1){2.0}}
%\put(8.9,0.75) {\line(1,0){2.2}}
%\put(11.1,0.75) {\vector(0,1){1.25}}
%\put(11.1,2.0) {\line(0,1){0.75}}
%\put(11.1,2.75) {\vector(-1,0){2.1}}
%\put(9.0,2.75) {\line(-1,0){0.1}}
%\end{picture}
%}
\newpage
%{\center

\begin{figure}[h]
\includegraphics[width=\textwidth]{flowchart8}
\caption{Time dependent coupling to plasma \ac{CFD} code}
\label{flow03}
\end{figure}
%{\Huge "time-dependent coupling", explicit}

%\vspace{2truecm}

%\unitlength1 true cm
%\begin{picture}(16,10) \thicklines
%
%\put(7.0,7.75){\framebox(2,1.5)}
%\put(7.2,8.75){\makebox(0,0)[tl]{LINDA}}
%
%\put(7.0,8.5) {\line(-1,0){4.5}}
%\put(2.5,8.5) {\vector(0,-1){1.5}}
%
%\put(9.0,8.5) {\line(1,0){4.5}}
%\put(13.5,8.5) {\vector(0,-1){5.5}}
%
%\put(1.5,5.5){\framebox(2,1.5)}
%\put(1.7,6.5){\makebox(0,0)[tl]{INTER0}}
%
%\put(2.5,5.5) {\vector(0,-1){1.5}}
%
%\put(0.0,0.0){\framebox(5,4.0)}
%\put(0.2,3.5){\makebox(0,0)[tl]{EIRENE}}
%\put(0.2,3.0){\makebox(0,0)[tl]{Source 1 in $\Delta t$}}
%\put(0.0,2.0) {\line(1,0){5.0}}
%\put(0.2,1.5){\makebox(0,0)[tl]{Source 2 at t=0,}}
%\put(0.2,1.0){\makebox(0,0)[tl]{from previous $\Delta t$}}
%
%\put(11.0,1.0){\framebox(5,2.0)}
%\put(11.7,2.5){\makebox(0,0)[tl]{BRAAMS}}
%\put(11.7,2.0){\makebox(0,0)[tl]{1 timestep}}
%
%\put(7.0,2.25){\framebox(2,1.5)}
%\put(7.2,3.25){\makebox(0,0)[tl]{INTER1}}
%
%\put(7.0,0.25){\framebox(2,1.5)}
%\put(7.2,1.25){\makebox(0,0)[tl]{INTER2}}
%
%\put(9.0,3.65) {\vector(-1,0){4.0}}
%\put(5.0,3.65) {\line(-1,0){0.1}}
%\put(4.9,3.65) {\line(0,-1){3.3}}
%\put(4.9,0.35) {\vector(1,0){2.1}}
%\put(7.0,0.35) {\line(1,0){4.1}}
%\put(11.1,0.35) {\vector(0,1){0.65}}
%\put(11.1,1.0) {\line(0,1){2.65}}
%\put(11.1,3.65) {\vector(-1,0){2.1}}
%\end{picture}
%}

%Source1 + Source2 : Stratified Source Sampling

%\medskip

%Source 1 : homogeneous in $\Delta t$

%\medskip

%Source 2 : bootstrapping from empirical distribution at t=0

\newpage

\begin{figure}{h}
\includegraphics[height=20cm]{converg}
\caption{Convergence of
B2/EIRENE, measured by B2-residuals. Method: implicit coupling,
method of ``saturated residuals", ref.\cite{kn:Maddison94}}
\label{converg}
\end{figure}

\subsection{correlated sampling}
A generic feature of coupled B2-EIRENE runs seems to be the loss of exponential decaying residuals (errors in B2 balances) down to machine precision,
which, in stand alone B2 runs, is a signature of having achieved convergence.
When coupled to EIRENE, at least with uncorrelated EIRENE trajectories between subsequent EIRENE calls (e.g.: between time steps
in the explicit coupling code), residuals only saturate. The level of the saturated residuals then scales only with the square root
of the number of EIRENE trajectories used per EIRENE cycle. Other convergence criteria then have to be defined, see \cite{kn:Maddison94}
for an early discussion of this still unsettled subject. The typical behaviour of residuals in coupled B2-EIRENE runs
(``saturated residuals'') is depicted in \cref{converg}.

Note:  strong (strict) convergence (exponentially decaying residuals down to machine precision)
can be achieved even in the presence of Monte Carlo noise, as is regularly
seen in simplified and dedicated test cases, and, actually, in the very first applications of EIRENE in coupled
plasma-neutral transport problems in the late eighties of the last century \cite{kn:Gerhauser88a}, there to \ac{TEXTOR} limiter configurations.
This, however, seems to require the activation of the ``correlated sampling" options of EIRENE
in order to provide positive statistical correlation between successive time-steps (or iterations).
However, except for strongly reduced (simplified) applications presently the level of positive correlation implemented via the EIRENE flag NLCRR (input block 1) seems
to be too low  to provide strong convergence on a regular basis and for full scale problems (e.g. the detached ITER divertor).
In these latter, more complex cases the entropy producing collisions and surface interactions lead to a rapid de-correlation
of trajectories, leading to a low level of correlation only, while, with the same algorithm, that level of correlation
can be very high in geometrically and/or physically very simple cases.
\\
\medskip

For this consider two plasma solutions obtained with a \ac{CFD} plasma code coupled iteratively to EIRENE,
at neighboring time-steps $1$ and $2$.
The incremental change of the plasma solution due to the change in the Monte Carlo reaction term in the plasma flow equations is the
difference between Monte Carlo source rates evaluated at time $1$ and time $2$.


In more general terms a Monte Carlo estimate  of a difference $\Delta E$ between two different calculations $MC_1$ and $MC_2$ (due to either code version differences, due to a different physical incremental effects or due to iterative procedures with \ac{CFD} solvers) is:
\begin{align}\label{cov}
\Delta E &= E(MC_1) - E(MC_2)  \nonumber \\
\sigma^2(\Delta E) &= \sigma^2(E(MC_1)) + \sigma^2(E(MC_2)) - 2 cov(E(MC_1),E(MC_2))
\end{align}
(This holds even for any statistical estimate of an incremental effect, independent of Monte Carlo procedures).
Here $E(MC)$ is the Monte Carlo estimate of the expectation value (i.e. it is the Monte Carlo solution to the kinetic transport equation),
and $\sigma^2$ is the corresponding statistical variance.
It is clear from \cref{cov} that small effects $\Delta E$ (such as those due to slightly different code versions) can only by detected when the statistical correlation quantified by the covariance of the two estimators $cov$, between the two Monte Carlo runs, is made large (EIRENE option: NLCRR).

\newpage

\section{non-linear effects: neutral-neutral collisions}  \label{sec1.8}
Neutral-neutral collisions make the Boltzmann collision integral
non-linear (bi-linear). More generally, inclusion of any binary
collisions between two test-particle species $i_0, j_0$, elastic
or not, have this effect. Important examples are, of course,
elastic self collisions $(i_0,i_0)_{el}$, e.g., $H_2-H_2$
collisions leading to viscous effects in the gas, elastic cross
collisions $(i_0,j_0)_{el}$, such as $H-H_2$ collisions, often
leading to a heating of the molecular gas due to hot charge
exchange atoms. Furthermore there may be relevant inelastic (chemical) binary
particle collisions $(i_0,j_0)_{inel}$, such as $(H_2-H^+_2
\rightarrow H+H_3^+)$ which is an important channel in hydrogen
plasma chemistry. Finally, binary test particle interactions occur in case of radiation transport, 
for example $(h\nu + H(1s) \rightarrow H(n=2))$, the
absorption of resonance Lyman-alpha photons. In the intermediate regime between optically thin and (black body) LTE
conditions the photons of a line from the discrete spectrum $h\nu$ become ``test particles'', 
which interact with other test particles, here with the ground state H atom, via line absorption.
In large dense divertor plasmas the ratio $H(1s)$ to
$H(n=2)$ may have to be calculated consistently with the radiation
field (a photon Boltzmann equation), competing electron impact
collisions and neutral atom transport.

Two methods of solution for this non-linear problem have been
tested in EIRENE, 1) a parametrization of the single particle
distribution functions involved in non-linear collision terms and
an iterative scheme (see next subsection: \ac{BGK} model collision
terms), and 2) a \ac{DMCS} technique
based upon a swarm simulation, an operator splitting (``Streaming"
and ``Collisions") and a random simulation of binary collisions in
the swarm. Currently only the first of these two options is still
being used. Despite its approximate character, it proved
to be much more effective than the full \ac{DMCS} method, for the cases
relevant to fusion research applications studied so far. A simple
estimate for the relevance of neutral-neutral gas collisions is
given by the so called Knudsen number $Kn$:
\begin{equation}
Kn = \frac{\lambda_{el}}{L} = \frac{1}{\sqrt{2} \sigma_{el} n_N L}
\end{equation}
Here $\lambda_{el}$ is the mean free path for elastic neutral
neutral collisions, $\sigma_{el}$ is the typical collision cross section (\si{\square\centi\metre}),
$n_N$ the gas density (\si{\per\cubic\centi\metre}) and $L$ is a characteristic dimension
(\si{\centi\metre}). Fixing the hard sphere molecular diameter at $d=4 \cdot
10^{-8}$ \si{\centi\metre} results in a typical cross section $\sigma_{el}$ of $(1-2) \cdot
10^{-15}$ \si{\square\centi\metre} to be used in this Knudsen number.  For a gas temperature  of
300 K and a pressure of 1 Pa ($n_N=2.414 \times 10^{14} \si{\per\cubic\centi\metre} )$ we then obtain:
\begin{equation}
\lambda_{el} (\si{\centi\metre}) = 1/(2.414\times 10^{14}\times 2\times 10^{-15}) \approx 2 ~\si{\centi\metre}
\end{equation}
In the \ac{BGK} approximation to the Boltzmann collision integral (see below, \cref{sec1-bgk}) only an ``effective'' mean free path
appears, because the collision rate $\nu_{BGK}$ is an effective (relaxation) rate only. When the common assumption of an velocity independent
collision rate is made, then this rate is
chosen such that the \ac{BGK} model fixes just
one single transport coefficient to a desired value in the continuum limit. E.g. either a certain viscosity, or  binary diffusion, or heat conduction, or any one
other selected transport coefficient in the continuum limit can be matched, when the Chapman-Enskog procedure is applied to derive continuum
equations (Navier Stokes equations) from the Boltzmann equation with that particular \ac{BGK} binary collision model term.

One default set of \ac{BGK} collision parameters used in EIRENE is based upon
the concept of ``diffusion volumes'' to provide (weakly temperature dependent) binary diffusion coefficients $D_{12}$ (for cross collisions, e.g. for He + H$_2$) or a viscosity coefficient $\mu$ (for self collisions) in the gas (see \cite{kn:Reiter94}). This choice results in a \ac{BGK} collision rate coefficient $k(T)= k_0 T^{0.25}$  with the constant $k_0$ derived
from such continuum limit considerations and the (experimentally determined) ``diffusion volumes''. Within this concept then a \ac{BGK} - mean free path \cite{kn:Kotov-Dis} follows as:
\begin{align}
\lambda_{BGK} (\si{\centi\metre}) &= \frac{\tilde{v}}{n \langle \sigma v\rangle} =  \frac{\sqrt{8T/\pi m}} {n k_0 T^{0.25}} \nonumber \\
                   &= L_0 \times T[\si{\electronvolt}]^{0.25} / n[10^{14} \si{\per\cubic\centi\metre}]
\end{align}
where the proportionality factor $L_0$ (in \si{\centi\metre}) is typically in the range between 7 and 20 ($L_0 = 20$\,\si{\centi\metre} for D--D self collisions, $L_0 = 9.75$\,\si{\centi\metre} for D$_2$--D$_2$ self collisions and $L_0 = 8.7$\,\si{\centi\metre} for D$_2$--He cross collisions). $L_0$ is the \ac{BGK} mean free path (in \si{\centi\metre}) in a gas
at $T = 1$ \si{\electronvolt} and density $n =1 \times 10^{14}$ \si{\per\cubic\centi\metre}.

For the same conditions as above: $T=300$ K ($T=0.026$ \si{\electronvolt}) and a pressure of 1 Pa ($n_N=2.414 \times 10^{14} \si{\per\cubic\centi\metre} )$, one then finds $\lambda_{BGK}\approx 1.6 ~\si{\centi\metre}$
in a $D_2$ gas.



For $Kn\ll 1$
a neutral fluid model (diffusion approximation, Euler, or Navier Stokes, etc.) is more appropriate,
while $Kn \gg 1$ leads to the (linear) free molecular flow regime, in which
only the linear collisions terms (collisions of neutrals with surfaces and a
given ``plasma'' host medium) are retained.
\subsection{\ac{BGK} approximation}\label{sec1-bgk}
currently under development. Please contact us for details and
status.

See flowchart 9.1 (\cref{flow6}).

Further details can be found in Ref. \cite{kn:Reiter96}.
\bigskip

The \ac{BGK} approximation used in the EIRENE code to model neutral-neutral collisions,
of both monatomic and multiple species gases,
and its implementation via successive approximation of the non-linear Boltzmann equation,
is a special case of the classical
``test particle method" for solving the Boltzmann equation, as opposed to the ``direct Monte Carlo
Simulation" method. The former was introduced in rarefied gas dynamics by J.K. Haviland, in 1961, see:
J.K. Haviland, ``The solution of Two Molecular Flow Problems by the Monte Carlo Method",
in: Methods in Computational Physics (Ed. Alder), Vol. 4, 1965, Academic Press.
An early review and comparison of different Monte Carlo methods in kinetic theory, also comparing this particular
``test particle method" with the \ac{DMCS} scheme, is given by
the report: N.A. Derzko, ``Review of Monte Carlo Methods in kinetic Theory", UTIAS Review No. 35 (April 1972)
University of Toronto, Institute of aerospace studies, April 1972.
(This review also includes a discussion of a third Monte Carlo scheme, the so called ``Integral evaluation method" of
A. Nordsieck and B.L. Hicks.)
\\
In rarefied gas dynamics problems for fusion reactor design, in particular for the ITER divertor design, the Haviland method,
but applied to a Boltzmann equation with a (non-linear) \ac{BGK} collision integral,
was first developed for the EIRENE code around 1995 (see Reiter et al. \cite{kn:Reiter96}),
and it is routinely applied in computational divertor design studies for ITER since
about 2004 by V. Kotov, A.K. Kukushkin, et al.

The method of ``successive approximations" employed in EIRENE is illustrated best when using the integral (Fredholm) form
of the Boltzmann equation......to be continued....


\subsubsection{A: Self collisions}


Consider first the elastic self collisions amongst particle species $A$.
The \ac{BGK} approximation to the Boltzmann collision integral reads:
\begin{equation}
\frac{\partial f_A(\vv)}{\partial t}_{BGK}=-\nu (f_A(\vv)-f^0(\vv))
\end{equation}
where
\begin{equation}
f^0(\vv)=\frac{n}{(2\pi mT)^{3/2}} \exp\left[-\frac{m(\vv-\vV)^2}{2T}\right]
\end{equation}
is a drifting Maxwellian distribution.

The parameters $n, \vV, T$ in $f^0$ are uniquely determined by conservation of particles, momentum and
energy in these elastic collisions.

When the collision rate $\nu$ is independent of velocity $v$, as it is the case in the default model described
above, then the parameters are simply given by the first velocity moments of $f_A$.

A key feature for these velocity independent collision frequencies is the simple closed form solution for $f_A(\vv,t)$
in a homogeneous medium, which translates into a particular simple realisation of this type of collisions in a Monte Carlo
procedure. Let $g(\vv)=f_A(\vv,t=0)$, then
\begin{equation}
f_A(\vv,t)=g(\vv)e^{-\nu t} +\left(1-e^{-\nu t}\right) f^0(\vv)
\end{equation}
where the parameters $(n, \vv, T)$ in $f^0$ can be calculated directly as moments of $g(\vv)$
(i.e. from the EIRENE default tallies PDEN, VDEN, EDEN.)

When a velocity dependent collision frequency is used, e.g. to capture
more than just one continuum limit transport coefficient correctly (or adjust to another Prantl number), then the solution of the
initial value problem for the distribution function $f_A$ is not as simple any more, and the
parameters $(n, \vv, T)$ in $f^0$ become fictitious local parameters related to the five moments of $f_A(\vv)$
weighted with $\nu(\vv)$:
\begin{equation}
\int \nu(\vv) M_i(\vv) f^0(\vv) d^3(v) =  \int \nu(\vv) M_i(\vv) f_A(\vv) d^3(v) \ \ \ i=1,\dots 5
\end{equation}
Here $M_i(\vv)$ are the five collisional invariants: $M_1 = 1$ (particle conservation),
$M_{2,3,4} = (v_x,v_y,v_z)$ (momentum conservation), and $M_5 = m/2 v^2$ (energy conservation).
Storage for the rates on the right hand side of these relations is foreseen on the EIRENE ``BGK'' default tallies,
and the Maxwellian integrals on the left hand side have to evaluated on a case by case basis, depending on the
choice of the velocity dependence in $\nu(\vv)$ made.


\subsubsection{B: Cross collisions}

to be written...

\subsubsection{B1: Cross collisions: electron or photon impact}

In case of photon absorption processes, or other processes with no momentum (velocity)
transfer, one virtual background species suffices:

For example, in case of resonance (Lyman) line absorption with excitation $H(1s) \rightarrow H(2)$,
the velocity distribution of H atoms does not change, and only a photon absorption rate
coefficient (additional to the electron impact ionisation rate coefficient) needs to be
provided, rather than a spectral parametrization of the Lyman radiation field, since only the density but not the
distribution function of the H atoms is affected by the photons.

From the photon point of view, a virtual parameterized H atom species (e.g. with flow velocity, temperature)
is needed, because via Doppler effects the absorption line shape in the photon transport equation
is very sensitive to the local atomic gas parameters   


\subsubsection{Implementation, ``virtual particle'' concept, iteration}

Self collisions between those particle species, which are
followed by EIRENE (i.e., between any two ``test-particle''-species)
can be included by using the iteration option NITER $> 1$  (input
block 1), and by adding a proper collision process to the collision
kernel (and its related total cross section $\Sigma_{nonlin}$ to the total cross
section $\Sigma$):
\begin{align}
C &= C_{lin} + C_{nonlin}\\
\Sigma &= \Sigma_{lin} + \Sigma_{nonlin}
\end{align}
Whereas $C_{lin}$ and $\Sigma_{lin}$ depend only upon parameters of the
background medium (``bulk particles'', usually: electrons and ions),
the kernels $C_{nonlin}$ and the rate coefficients $\Sigma_{nonlin}$ may
also depend upon parameters of the test particles (atoms, molecules,
test ions) themselves.

If other non-linearities are already present in the model, such as those due to
self-consistent feedback between the neutrals and the plasma state
(see above: \cref{sec1.7}), code-iterations between the test particle and bulk plasma parameters can
also be used to simultaneously iterate on the neutral-neutral collision effects. Then NITER=1 suffices, i.e. only one single EIRENE
cycle per plasma code step, to relax simultaneously
both the neutral-neutral interaction and the neutral-plasma interaction non-linearities.

In order to include such non-linear interactions in a case, a copy of
the test-particle species cards (input block 4a, 4b or 4c) is put into
the background species specification (input block 5), as a ``virtual'' species.
E.g., if one
wishes to include H - H elastic neutral-neutral collisions in a
simulation, then a (virtual) species of H particles with a drifting Maxwellian distribution has to be
specified as one additional background particle species (``bulk ions with charge
state zero'', by abuse of language). During the
iterative
procedure, the parameters of the Maxwellian of this artificial background species
(e.g. density, temperature, drift velocities in case of velocity independent collision rates)
are those that have been
obtained from the corresponding test particle tallies in the previous
iteration. One iteration step ends by overwriting these parameters
at the end of a cycle with the new values, in the post processing phase
of each iteration cycle (module MODBGK.f) and by preparing the new collision rate coefficients
for the next iteration cycle.



\subsection{Direct Simulation (\acs{DMCS}) of self-collisions}
The \ac{DMCS} method was implemented in EIRENE in the late eighties of the last century,
but then abandoned due to its limitations to physically rather simple cases (e.g. no feedback
with the consistent bulk plasma state, the often dominant non-linearity in fusion applications).
Currently it is not developed nor supported  any more. Contact us for details and
status. Further details can be found in Ref. \cite{kn:Behringer}

\section{Radiation transport, photon gas simulations}
% dissertation Veach, 1997
The extensions of EIRENE towards radiative transfer problems, and also the material in this
present short description, is largely based on the monograph by Oxenius: \cite{kn:Oxenius}.

The theory of light and its interaction with matter is one of the
best that physics can offer, and can predict virtually every
observed phenomenon with great accuracy [Feynman, 1985, QED: The
strange Theory of Light and Matter, Princeton University Press,
Princeton]. In our simulations of radiation transport with the
EIRENE code we assume a geometric optics model (i.e., essentially,
``particle theory of light''), which is sufficient for all our
purposes. Light is emitted, scattered, reflected and absorbed, but moves on
straight lines between these events. I.e., a continuously varying
index of refraction is not allowed, and we ignore most properties
of light that depend on a wave or quantum model for their
explanation (e.g. diffraction, or interference). However, the
linear wave-optics phenomenon of polarization and the linear
quantum optics phenomenon of fluorescence [absorption of a certain
wavelength (line), and re-emission at another wavelength
(line)] can conveniently be added to our linear geometric optics model.\\
% Oxenius
We are concerned with the kinetic theory of interacting particles
and photons, where the behaviour of each species is governed by a
kinetic equation: the kinetic equation for photons being nothing
else but the well known equation of radiation transfer, i.e., the
``strangely normalized photon-Boltzmann equation''. In contrast to
plasma radiation (e.g., Bremsstrahlung) the radiation processes
studied here are due to single, uncorrelated particles such that
emission and absorption of the plasma is obtained by simple
summation of all contributions of all individual particles. We
therefore consider all particles as uncorrelated and in well
defined one-particle states.

For continuum radiation (bound-free or free-free) the frequency of
the photon in the atoms rest frame may be replaced by the
frequency in the lab. frame, because all cross sections are slowly
varying functions of frequency. By contrast, for line-radiation
due to bound-bound transitions, the Doppler effect must be taken
into account in all quantities that vary rapidly as function of
frequency, such as cross sections, line-profile, or redistribution
functions
%Feynman, Pomraining, Oxenius, Uhlenbusch,...

The key atomic parameters for radiation transport are
the Einstein (or Einstein-Milne) coefficients A$_{mn}$, B$_{mn}$
and B$_{nm}$ of spontaneous emission, induced (stimulated) emission
and absorption, respectively. We will discuss them here first for
a gas of stationary atoms and include Doppler- and other line
broadening effects later.

For a thermal radiation field (temperature T) in a gas of atoms
with density $n_1$ (lower state) and $n_2$ (upper state) the rate
equation expressing balance between emission and absorption reads:
\[n_1 B_{12} R_\nu(T)= n_2[A_{21}+B_{21} R_\nu(T)]\].

$R_\nu(T)$ is a specific function characterizing the wavelength and
temperature dependence of the
thermal radiation field (i.e., the Planck-function). For example
the specific intensity $I_\nu$
(energy\-/area\-/time\-/solid-angle\-/frequency-interval) is often used, or the
(angle)-integrated intensity $J_\nu = \int I_\nu d\Omega$, or the
specific energy density $u_\nu = 1/c J_\nu$
(energy\-/volume\-/frequency-interval), etc. Also the wavelength
dependence may be given either using frequency $\nu$ (Hz), $\omega$
(\si{\radian\per\second}), wavelength $\lambda$ or energy $E = h \nu=\hbar \omega$.

Depending upon all these choices the numerical factors in the
definitions of the $B_{nm}$ coefficients (sometimes even in the
$A_{nm})$ differ in the literature. For our purpose (Monte Carlo
solution of the full radiation transfer equation) we need to
define monochromatic mean free pathes for photons,
rates and rate coefficients
and cross sections for the individual processes.

In order to take advantage of the analogy between radiation and particle
transport equations as much as possible, we will choose the numerical constants
in the Einstein coefficients such that these
parameters have the same meanings and even the same units in both cases of
radiation and particle transport.
I.e., we will use the
``absorption coefficient'' (better: ``extinction coefficient'',
if scattering is allowed)
$\alpha_{\nu} = n \sigma \equiv \Sigma_t$,
(invers monocromatic mean free path), which directly corresponds to
the ``total macroscopic cross section'' (dimension: 1/length) defined earlier for
particle transport problems.


% d.h. definitionen der b_nm muessen mit den formeln fuer sigma, etc...
% consistent gemacht werden!!!!

%weiter r.j.rutten, (rad.transf. in stellar atmospheres, hat gute
%diskussion und saubere definitionen. insb. a_nm explizit, (4 pi
%factor bereits drin, gibts auch anders, bereits bei a_mn,
% bei uns: so wie dort.
% aber vorsicht: feynman
%hat I(omega) verwendet, deshalb einen faktor 2pi anders zwischen
%a12 und b12 als rutten, rybiki,oxenius,wir.
%frage: wir verwenden energie, nicht frequenz. ist da alles richtig?
%uhlenbusch,einstein,chandrasekhar,.... verwenden u(nu), also
% um faktor 4pi/c andere
%relation zwischen a12 und b12.
\subsection{Line shape options}
Emission of photons from a particular line is simulated by sampling
firstly the location of emission, and secondly the frequency and direction of emission.
Spatial distribution of photon emission is treated in EIRENE in exactly the same way as
for any other particle type (neutrals, ions), and actually using the same modules of code for that.
Emission from the transition $i\rightarrow k$ from an upper state $i=nl$ of an atom is treated, by abuse  of language, as ``volume recombination
source": The position is sampled from the spatial distribution $A_{ik}P^{nl}(\vec{r})$, which, because the Einstein
emission coefficient $A_{ik}$ is constant, is the same as the spatial distribution $P^{nl}(\vec{r})$ of upper state particle density.
.....
Given the birth point $\vec{r}$ of the photon, and hence  all local parameters of the host medium at this point,
the frequency $\nu$ (in fact: the energy $E=h\nu$) is sampled from the normalized line shape function $\Phi(\nu)$,
....
The following line shape options are currently available:
\begin{itemize}
\item{ISHAPE=0} $\delta$-profile: mono-energetic emission (not yet available for absorption)
\item{ISHAPE=1} pure Doppler profile
\item{ISHAPE=2} Lorentz profile, truncated at zero
\item{ISHAPE=3} Doppler convoluted Lorentz  profile, i.e. Voight profile
\item{ISHAPE=4} LORVDW: convoluted Lorentz and van der Waals pressure broadened profile
\item{ISHAPE=5} Doppler broadened LORVDW profile
\item{ISHAPE=6} normal Zeeman triplet, $\delta$-profile in each component
\item{ISHAPE=7} Doppler broadened Zeeman-$\delta$-profile in each component
\item{ISHAPE=8} Doppler broadened Zeeman-Stark-profile, 10 Lorentzian model, Rosato 2006
\end{itemize}



\newpage
\section{Charged Particle Transport}
EIRENE also follows charged particles since EIRENE$_{1987}$. The particle tracing is
carried out in Subroutine FOLION.
We refer to this ``default" ion transport model as the \emph{``minimal trace ion model"}, as opposed to
refinements added only later, starting with Seebacher et al., (2012) \cite{kn:Seebacher}.
The differences of the module FOLION to
the module FOLNEUT for neutral particles (and photons) are:
\begin{itemize}
\item Motion is not along straight lines between events, but instead
determined by drifts and forces due to the electro-magnetic fields.
\item Time stepping is not of ``Poisson type" (exponentially distributed time steps), but with constant increments $\Delta t$,
because the equation of motion has to be solved numerically, and this introduces already a discrete time-stepping anyway.
\item after each time-step a Coulomb collision kernel is modelled,
i.e., a diffusive step in velocity space.
\item after each time-step a diffusion-step is modelled,
i.e., a diffusive step in physical space (not yet fully written).
\item ionized test particles which hit a surface see an electrostatic sheath potential,
by which they are either reflected, slowed down or accelerated,
depending on the sign of their charge relative to the sheath
potential.
\end{itemize}
\subsection{Orbit integration}
Orbit integration in a magnetized plasma is carried out to various orders
of the small parameter $\delta = \rho/L$, ($\rho$:
thermal gyro-radius, $L$ spatial length scale, e.g. gradient length of B-field, etc.)
\subsubsection{Particles following the B-field}
Simple, analytical, orbit integration, zeroth order in $\delta$.
\paragraph{A: no electric field, no drifts}\hfill
To zeroth order in the smallness parameter $\delta=\rho/L$ , i.e. for $\rho \simeq 0$,
charged particle orbits are the magnetic field lines themselves.
The magnetic field unit vector $\textbf{b}=\textbf{B}/|\textbf{B}|$ is provided as input
vector in each cell of the computational
grid, see input block 5 (\cref{sec2.5}). By default the magnetic field vector
is taken to be constant within a grid cell, as are also the plasma (background) input tallies.

With the  ``FEM" interpolation option (EIRENE$_{2006}$ and younger) iso-parametric elements can be used
for interpolating input tallies in various grid options, and hence providing a magnetic field $\textbf{B}$ which continuously varies
inside the grid cell (and then also providing non-zero derivatives such as grad|$\textbf{B}$| or B field curvature).

In 2D cases the interpolation scheme currently ignores the third coordinate,
even if it is curvilinear (e.g., NLTRA option, toroidal effects in a 2D setting).
As a consequence toroidal curvature effects are not yet properly
seen by charged particles in that case, and the corresponding toroidal curvature drift had been added
separately in \cite{kn:Seebacher}.

Default is $\textbf{b}=(0.,0.,1.)$ in each cell.
The velocity vector of each test-ion followed by EIRENE is projected onto $\textbf{b}$
in each grid cell along the trajectory. The velocity used for pushing the particle along
its trajectory is only the parallel component $\vv_\|=(\vv \cdot \vb) \vb$
\paragraph{B: \textbf{E} $\times$ \textbf{B} drift only}\hfill
Only the \textbf{E} $\times$ \textbf{B} drift is included.
This is still zeroth order in $\delta$.
By modifying the electric field
in each cell into one effective \textbf{E$^*$} field, some features of  guiding centre drifts
(namely: the grad B drift) can be already captured at this level.
\subsubsection{Guiding centre approximation}
To first order in $\delta$ a number of so called guiding centre drifts arise. Their inclusion
into EIRENE (module FOLION) is a later code extension, and not fully operational yet.
Once this option is fully implemented, guiding centre currents will be derivable from EIRENE runs
as volume tallies. These guiding centre currents, supplemented by the (divergence free)
magnetisation currents, are one important component of the plasma currents.
\subsection{Coulomb Collision Models}
In this section conventional notation for masses of colliding charged test particles labeled $a$ and field (background) particles $b$
is used:   $m_a, m_b$ for the masses,
for atomic mass numbers
$\mu_a, \mu_b$ and for reduced
masses $\mu_{a,b}$ are used. Charge states are $Z_a, Z_b$. Particles of type $b$ are background particles,
with a given velocity distribution $f(\vec v_b)$, whereas particles of type $a$ are individual test particles with a specific
velocity $\vec v_a$.


Four (velocity dependent) collision relaxation rates (inverse time) arise: The slowing down rate
$\nu_s^{a,b}(v_a)$, the deflection (or transverse diffusion) rate $\nu_\perp^{a,b}(v_a)$, dispersion (or parallel diffusion) $\nu_\parallel^{a,b}(v_a)$ and the energy loss rate $\nu_\varepsilon^{a,b}(v_a)$

Due to energy conservation in the binary Coulomb collisions one has \cite{kn:NRL}:
\begin{equation}\label{eps1}
\nu_\varepsilon^{a,b}(v_a) = 2 \nu_s^{a,b}(v_a) - \nu_\perp^{a,b}(v_a) - \nu_\parallel^{a,b}(v_a).
\end{equation}
I.e.: the energy transfer rate $\nu_\varepsilon^{a,b}$ is given via the ``slowing down rate" $\nu_s^{a,b}$, the transverse diffusion
rate $\nu_\perp^b$ and the parallel diffusion rate
$\nu_\parallel^b$, for any velocity $\vec v_a$ and any distribution $f_b(\vec v_b)$.

Explicit results when $f(\vec v_b)$ is a Maxwellian distribution are given in most textbooks or lecture notes
on plasma physics since many decades. In general: unless otherwise stated, notations here are as in \cite{kn:NRL}.
The resulting rates (or: rate coefficients, if divided by target density $n_b$) have the same functional dependencies
on test particle velocity (or energy) and background temperature as the ``Beam-Maxwellian rates" (or rate coefficients)
of format H.3, H.4, H.5 for in-elastic collisions described in section 4 (atomic and molecular data input, \cref{sec2.4})

\subsubsection{Simple Coulomb Collision Models}
Strongly simplified Coulomb collision models are usually sufficient for molecular ions,
because of the molecular (radical) ion's typically very short lifetime until further fragmentation (dissociative excitation,
dissociative recombination, etc...) to atomic products. In order to roughly account for thermalization effects
and background plasma
cooling or momentum transfer,  the energy, parallel and the perpendicular velocity of such ions can be
modified ``continuously" during particle orbit integration, hence approximating the velocity space
diffusion (Fokker Planck equation) by mean values of energy and momentum of the test particle.
There are two classes of simple collision models: deterministic and stochastic.
\paragraph{A: Deterministic relaxation in velocity space}
The energy and/or velocity of charged test particles are modified in a fully deterministic way along the
particle trajectory. I.e.: along with orbit integration also energy, or certain velocity components
are evolved according to ordinary (non-stochastic) differential equations, using appropriate relaxation times.
\subparagraph{A1: Only energy relaxation}\hfill \\
The simplest Coulomb Collision model is an averaged \acs{BGK}-type
relaxation of energy of test particles towards the mean energy of the Maxwellian
background ions. In the earliest versions of the EIRENE code (mid
eighties) it was originally based upon the hydrocarbon ion
equilibration model given by Langer, \cite{kn:Langer} and it is a slight
generalization thereof. The equation for the energy $\varepsilon_a$ of
a test-ion $a$ in a field of background ions $b$ (field particles) with temperature
$T_b$ is given as:
\begin{equation}\label{eq.relax1}
\frac{d\varepsilon_a}{dt}=-\nu_\varepsilon^{a,b} \cdot
\varepsilon_a  \quad , \quad  \varepsilon_a(t=0)=\varepsilon_a^0
\end{equation}
In general this is a non-linear equation because the energy exchange collision
rate (energy transfer rate) $\nu_\varepsilon^{a,b}$ depends on
$\varepsilon_a$ itself. The full expression for $\nu_\varepsilon^{a,b}(v_a, T_b)$ is given below, it involves
the Maxwell integral $\Psi$ and its derivative $\Psi'$, see below.



\vspace{1cm}
The \emph{limit of small test particle velocities} (as compared to thermal velocities of the
field particles $b$), i.e.:
\begin{equation}\label{xab}
v_a\ll \sqrt{(2T_b/m_b)},\quad  \mbox{i.e.} \quad  x^{a/b} = \mu_bv_a^2/2kT_b \ll 1,
\end{equation}
as often particularly
appropriate for molecular ions $a$,
In this limit one has for the basic ``Spitzer
rates" (loc. cit.):
\begin{align}
\nu_s^{a,b} & = n_b Z_a^2 Z_b^2\lambda_{a,b}\cdot
6.8 \times 10^{-8}~\frac{\mu_b^{1/2}}{\mu_a} \left(1+\frac{\mu_b}{\mu_a}\right)T_b^{-3/2}\\
\nu_\perp^{a,b} & = n_b Z_a^2 Z_b^2\lambda_{a,b}\cdot 1.4 \times 10^{-7}~\frac{\mu_b^{1/2}}{\mu_a} T_b^{-1/2}\varepsilon^{-1}  \\
\nu_\parallel^{a,b} & = n_b Z_a^2 Z_b^2\lambda_{a,b}\cdot 6.8 \times 10^{-8}~\frac{\mu_b^{1/2}}{\mu_a} T_b^{-1/2}\varepsilon^{-1}
\end{align}
and hence, using Equation \cref{eps1}:
\begin{align}
\nu_\varepsilon^{a,b} & = n_b Z_a^2 Z_b^2\lambda_{a,b}\cdot
6.8 \times 10^{-8}~\frac{\mu_b^{1/2}}{\mu_a}\frac{1}{T_b^{1/2}}
\left[2 \frac{1}{T_b}\left(1+\frac{\mu_b}{\mu_a}\right)-
2\frac1\varepsilon_a-\frac1\varepsilon_a\right]\\
                      & = \nu^c_\varepsilon +\nu^v_\varepsilon
\end{align}
with the ``constant" and ``variable" contribution $\nu^c_\varepsilon$ and $\nu^v_\varepsilon$, respectively, defined as
\begin{align}
\nu^c_\varepsilon &= 2 n_b Z_a^2 Z_b^2\lambda_{a,b}\cdot
6.8 \times 10^{-8}~\frac{\mu_b^{1/2}}{\mu_a}\frac{1}{T_b^{3/2}}
\left[1+\frac{\mu_b}{\mu_a}\right] =
\tilde{\nu}_\varepsilon\left[1+\frac{\mu_b}{\mu_a}\right] \\
\nu^v_\varepsilon &= 2 n_b Z_a^2 Z_b^2\lambda_{a,b}\cdot 6.8 \times 10^{-8}~\frac{\mu_b^{1/2}}{\mu_a}\frac{1}{T_b^{3/2}} \left[-\frac{3T_b/2}{\varepsilon_a}\right]=\tilde{\nu}_\varepsilon\left[-\frac{3T_b/2}{\varepsilon_a}\right].
\end{align}
Using this decomposition of $\nu_\varepsilon^{a,b}$ in equation
\cref{eq.relax1} we find:
\begin{equation}\label{eq.relax2}
\frac{d\varepsilon_a}{dt}=-\tilde{\nu}_\varepsilon\left[1+\frac{\mu_b}{\mu_a}\right] \cdot
\varepsilon_a + \tilde{\nu}_\varepsilon\cdot \frac{3T_b}{2}
\end{equation}

Hence, in this small test particle velocity limit of $\varepsilon_a \ll 3/2T_b$,
and introducing the mass correction factor $\delta_{a,b}=1+\mu_b/\mu_a$ the solution of this (now: linear) ordinary differential
equation exists in closed form:
\begin{equation}\label{eq.solut1}
\varepsilon_a(t)=\varepsilon_a^0 \cdot
\exp(-\tilde{\nu}_\varepsilon\delta_{a,b}  t)+
\frac{3T_b}{2\delta_{a,b}}\left[1-\exp(-\tilde{\nu}_\varepsilon \delta_{a,b}t)\right]
\end{equation}

\textbf{Example: hydrocarbon ions \cite{kn:Langer}}

With $\mu_b=1$ for protons, $\mu_a=16$ for CH$_4^+$, the Coulomb logarithm $\lambda_{a,b}\approx 10$,
hence $2 (\lambda_{a,b}/\mu_a) \delta_{a,b} \approx 1.3$ and
\begin{equation}
{\nu}_\varepsilon = \tilde{\nu}_\varepsilon \delta_{a,b}\approx
8.8 \times 10^{-8}~ n_b~T_b^{-3/2}
\nonumber
\end{equation}
($n_b, T_b$ in \si{\per\cubic\centi\metre} and \si{\electronvolt}, resp.)
which is the solution Eq. (11) given in \cite{kn:Langer} for thermalization
of hydrocarbon ions in a Maxwellian hydrogen plasma. The relaxation is towards an energy
$\varepsilon_p = 3/2 T_b/(1+1/16)$, which is (very slightly) lower only than the mean energy of the target Maxwellian, due to
the particular velocity dependence in the energy relaxation rate in this small test particle velocity limit.

And this is also exactly the approximation which was also used in EIRENE
since about 1987 for all test ions traveling in a bath of background ions.

\textcolor[rgb]{1.00,0.00,0.00}{To be done: in case of drifting Maxwellian background: first transform into plasma frame, such that stationary Maxw. background,
then relax test particle energy in this frame, then transform back...
To be done: isotropization: currently the ratio $v_{\perp B}$ to $v_{\parallel B}$ relative to the B-field  is preserved.}

\vspace{1cm}
The \emph{full expression for the energy relaxation frequency }$\nu_\varepsilon^{a,b}$, without the approximation of slow particles $a$, i.e. for all values of $x^{a/b}$ is:
\begin{equation}\label{eps2}
\nu_\varepsilon^{a,b}(x^{a/b}) =2 \left[\mu_a/\mu_b \Psi(x^{a/b}) - \Psi'(x^{a/b})\right] \nu_0^{a/b}
\end{equation}
where $x^{a/b}$ has been defined above, relation \cref{xab}
\begin{equation}
\nu_0^{a,b}  =  n_b Z_a^2 Z_b^2\lambda_{a,b}\cdot 4 \pi  /(\mu_a^2~  v_a^3)
\end{equation}
 and the Maxwell integral $\Psi(x)$:
\begin{align}
\Psi(x) &= \frac{2}{\sqrt{\pi}}\int_0^x dy~e^{-y} \sqrt{y} &&= \mathrm{erf}(\sqrt{x})-\frac{2}{\sqrt{\pi}}e^{-x}\sqrt{x}\\
\Psi'(x)&= d\Psi/dx &&= \frac{2}{\sqrt{\pi}}e^{-x}\sqrt{x}
\end{align}
With this more general expression for the collision frequency Equation \cref{eq.relax1}
can then only be integrated numerically  along the particle trajectory.
However, other features of the collision process (velocity spreading, isotropization)
then should be taken into account as well in order
to ensure proper equilibration of energy $\varepsilon_a$ of test particles
to $3/2 T_b$ of the fixed (stationary Maxwellian) background, as well as momentum conservation.
Without such further velocity space effects
a test-particle with energy $\varepsilon_a$ corresponds to the non-equilibrium situation
of a plane parallel (mono-energetic) flux, and energy exchange of this non-equilibrium flux
with a Maxwellian background
vanishes at a critical energy $\varepsilon'$ distinct from from $3/2 T_b$.
\cite[paragraph 19]{kn:Trubnikov}.
\subparagraph{A2: \ac{BGK} Coulomb-Collision operator}\hfill \\
(Diploma Thesis Felix Reimold....)
\subsubsection{The Fokker-Planck Collision Model}
\paragraph {A:}
A gyro averaged Fokker-Planck Coulomb Collision Model, which was originally
developed for a stand alone trace impurity ion transport code DORIS
\cite{kn:Doris,kn:Reiser98}. It was later accommodated as collision kernel in EIRENE.
Distinct from various other drift-kinetic collision operators published already since the early
nineties of the last century its range of applicability is rather limited.
In its present version this collision operator ignores the difference between
``real particle position" and the guiding center, hence, e.g. (classical) diffusion terms
due to finite Larmor radius effects are not taken into account. The averaging over gyro angle
limits this operator to
distributions of relative velocities which are symmetric around the B-field line, i.e. essentially a 1D, parallel to B, situation,
(but it allowed for distributions which may deviate from Maxwellians in this B-field direction,
e.g. to accommodate parallel to B thermal force effects).
See also: Thesis J. Seebacher, Univ. Innsbruck, 2009, and \cite{kn:Seebacher}.

\paragraph{B: binary Monte Carlo collision procedure}
A procedure similar to the ``Abe, Takizuka Coulomb collision algorithm" frequently used in PIC codes \cite{kn:Takizuka}
is employed. This procedure consists in random sampling  a bulk (background) collision partner, and then to carry out
an effective binary collision between these two particles. It hence can be viewed as Monte Carlo ``randomization"
of the algorithm A, in which the integration over background distribution is carried out deterministically.


\newpage
\section{Parallelization of EIRENE Code}
EIRENE is often coupled iteratively to finite volume or finite element \ac{CFD} codes, which provide the host medium (plasma background).
Then EIRENE is used as kinetic neutral particle transport module,
providing, often in iterative mode, the reaction part in the overall
``diffusion-advection-reaction problems. Non-linearly coupled fusion edge plasma flow
simulations fall into that category.

In contrast to the \ac{CFD} codes, which can be parallelized using domain decomposition, EIRENE uses a particle partitioning technique,
which is usually most suitable for Monte Carlo methods. Here the stochastic Monte Carlo trajectories are distributed on different processes,
whereas the domain is the same on all processes. If there would be no stratified source sampling also typically involved, then
this particle partitioning would provide a good load balancing automatically. Most 3D edge plasma applications using the EMC3-EIRENE
hybrid Lagrangian \acs{CFD}/Lagrangian kinetic Monte Carlo code are run in this (``embarrassingly parallel") single stratum mode of operation.

However, the 2D hybrid finite-volume/Lagrangian Monte Carlo code B2-EIRENE is heavily resorting on ``stratified sampling". This means the primary sources
of neutral particles are split by their physical or computational origin into ``sub-strata", and results are then linearly superimposed finally.
This is, amongst others, a variance reducing technique, but it also leads to load balancing issues in the Monte Carlo part, even with particle partitioning.



Different strategies have been implemented in the past
for assigning strata to processes:

problem specific matrix: PROCFORSTRA(NSTRA,0:NPRES-1),
to be filled in startup routine pedist.f

.... to be written....
\newpage

\section{EIRENE flow charts} \label{sec1.9}


\begin{figure}[h]
\includegraphics[width=14cm]{flowchart1.eps}
\caption{Structure:
EIRENE code}
\label{flow1}
\end{figure}

\newpage

\begin{figure}
\includegraphics[height=\textwidth]{flowchart2}
\caption{Flowchart:
Program EIRENE}
\label{flow2}
\end{figure}

\newpage

\begin{figure}
\includegraphics[height=22cm]{flowchart3}
\caption{Flowchart 2: Subroutine INPUT}
\label{flow3}
\end{figure}

\newpage
\begin{figure}
\includegraphics[height=22cm]{flowchart4}
\caption{Flowchart 6: Subroutine MCARLO}
\label{flow4}
\end{figure}

\newpage

\begin{figure}
\includegraphics[height=16cm]{flowchart5_new}
\caption{Flowchart 6.3:
Particle loop in subr. MCARLO.  Default volume averaged tallies are updated, ``on the fly" \emph{(per event)} in routines UPDATE,
problem specific tallies
for code-code interfacing (COPV(...), e.g. B2-EIRENE) in UPDCOP, specific tallies for non-linear \ac{BGK} operators (BGKV(...)) in UPDBGK.
Corresponding variances
are updated \emph{(per history)} after completion of each trajectory in STATS1, STATS1\_COP, STATS1\_BGK.
\textcolor[rgb]{1.00,0.00,0.00}{Addition/revision in 2012: Updating of linear functions of default tallies, \emph{per history}, in UPFCOP, e.g. for code coupling tally
COPV.
This makes redundant the specific statistical variance evaluation in STATS1\_COP}
}
\label{flow5}

\end{figure}
\newpage

\begin{figure}
\includegraphics[height=22cm]{flowchart6}
\caption{Flowchart
9.1: Iterative mode for neutral-neutral interactions}
\label{flow6}
\end{figure}
